\documentclass[paper=letter,fontsize=11pt,titlepage,openany,twoside,
numbers=noenddot,index=totoc]{scrbook}
%\usepackage[dvips]{graphicx}
\usepackage{etex}
\usepackage{graphics}
%%%%%%%%%%%%%%%%%%%%%%%%%%%%%%%%%%%%%%%%%%%%%%%%%%%% 
%%%%%%%%%%%%%%%%%%%%%%%%%%%%%%%%%%%%%%%%%%%%%%%%%%%%%%%%%%

\usepackage{import} %%%%%%
%\usepackage[pdflatex]{graphicx}
\usepackage{float}
\usepackage[intoc]{nomencl}
\usepackage{xfrac}
\usepackage{fancybox}
\usepackage{pifont}
\usepackage{phaistos}
\usepackage{cancel}
\usepackage{amssymb}
\usepackage{amsmath}
\usepackage{amsthm}
\usepackage{empheq}

\usepackage{xcolor}
\usepackage[english]{babel}
\usepackage{fancyhdr}
\usepackage{fancyvrb}
\usepackage{shadethm}
\usepackage{makeidx}
\usepackage[T1]{fontenc}
\usepackage{fouriernc,bm}
\usepackage[scaled=0.80]{beramono}
\usepackage{ragged2e}
\usepackage{paralist}
\usepackage{mdwlist}
\usepackage{multicol}
\usepackage[labelsep=none]{caption}
\usepackage{subfig}
\usepackage{tikz}
\usetikzlibrary{arrows}
\usetikzlibrary{patterns}
\usetikzlibrary{decorations}
\usetikzlibrary{intersections}
\usetikzlibrary{matrix}
\usetikzlibrary{snakes}
\usetikzlibrary{calc}
\usetikzlibrary{backgrounds}
\usepackage{picins}
\usepackage{framed}
\usepackage{nextpage}
\usepackage{colortbl}
\usepackage{listings}

\lstloadlanguages{Gnuplot}
\usepackage[pdfproducer={pdflatex},
            pdfauthor={Martin Veillette},pdftitle={Mathematical Methods in Physics},
            pdfsubject={Mathematical Methods in Physics},pdfkeywords={Math} {Physics},
            pdfdisplaydoctitle=true,bookmarks=true,pdfborder={0 0 0},colorlinks=true,linkcolor=cyan, urlcolor=blue,
            bookmarksnumbered=true,hyperindex,pdfpagelabels]{hyperref}
%\usepackage[colorlinks, hyperindex, plainpages=false, linkcolor=blue, urlcolor=blue, pdfpagelabels]{hyperref}
\usepackage{breakurl}

%%%%%%%%%%%%%%%%%%%%%%%%%%%%%%%%%%%%%%% for introduction
%%%%%%%%%%%%%%%%%%%%%%%%%%%%%%%%%%

\newtheorem{ex}{\bf Example}[section]
\newtheorem{optionaltopic}{\bf Optional Topic}[section]
\newtheorem{homework}{\bf Homework}[section]
\newtheorem{eg}{\bf Example}[section]
\newtheorem{selfcheck}{\bf Selfcheck}[section]
\newtheorem{dq}{\bf Discussion Section}[section]
\newtheorem{indentedblock}{\bf Discussion Section}[section]
\newtheorem{hwsection}{\bf Discussion Section}[section]

\newcommand{\hwendpart}{} % DONT KNOW
\newcommand{\answercheck}{} % DONT KNOW
\newcommand{\hwsoln}{} % DONT KNOW
\newcommand{\worked}[1]{#1} %DONT KNOW
\newcommand{\startdq}{} %DONT KNOW
\newcommand{\myvec}[1]{\mathbf{#1}} % Use /vc{} instead.
\newcommand{\myunit}[1]{\textup{#1}}
\newcommand{\kgunit}{\myunit{kg}}
\newcommand{\gunit}{\myunit{g}}
\newcommand{\munit}{\myunit{m}}
\newcommand{\sunit}{\myunit{s}}
\newcommand{\degunit}{\ensuremath{\mspace{0mu}^\circ}} % the \mspace{0mu} is a zero-width math space; need it so that, e.g., $2\times10^{-3}\degunit$ doesn't get interpreted as a double superscript
\newcommand{\degcunit}{\degunit\textup{C}}
\newcommand{\junit}{\myunit{J}}
\newcommand{\nunit}{\myunit{N}}
\newcommand{\kunit}{\myunit{K}}
\newcommand{\unitdot}{\!\cdot\!}
\newcommand{\Degrees}{\ensuremath{\mspace{0mu}^\circ}} % the \mspace{0mu} is a zero-width math space; need it so that, e.g., $2\times10^{-3}\degunit$ doesn't get interpreted as a double superscript
%-------------------- math in sans-serif text -------------------
% $...\zs{x}...$ -- italic, for use when I need symbol font mixed in
% $...\zu{x}...$ -- upright, e.g. for digits and things in sans
% See DeclareMathAlphabet near top, and p. 177 of LaTeX companion, for \ms{}.
% Note: don't put superscripts or subscripts inside these, or any math commands
% at all.
% They can all be used outside of math mode in examples or other sans-serif
% text, but won't necessarily produce the right results in serif text,
% since zs forces sans serif. So when cutting and pasting from example text
% into serif text, have to change zs's.
\newcommand{\zs}[1]{\ensuremath{\ms{#1}}}
\newcommand{\zu}[1]{\textup{#1}}
\newcommand{\zb}[1]{\textbf{#1}}

%%%%%%%%%%%%%%%%%%%%%%%%%%%%%%%%%%%%%%%%%%%%%%%%
%%%%%%%%%%%%%%%%%%%%%%%%%%%%%%%%%%%%%%


% Set the sans serif font to Helvetica (for chapter and section headings)
\renewcommand\sfdefault{phv}

\definecolor{captioncolor}{HTML}{005587}
\definecolor{linecolor}{HTML}{0074C8}
\definecolor{fillcolor}{cmyk}{0.1,0.05,0,0}
\definecolor{brickcolor}{HTML}{F0D8B2}
\definecolor{blockcolor}{HTML}{B6B6B6}
\definecolor{groundcolor}{HTML}{E4D8C5}
\definecolor{earthcolor}{HTML}{C5FFFF}
\definecolor{watercolor}{cmyk}{0.1,0.05,0,0}
\definecolor{codecolor}{HTML}{FFF7E0}
\definecolor{linenumcolor}{HTML}{8C8C8C}
%\definecolor{headercolor}{HTML}{F5F5F5}
\definecolor{headercolor}{HTML}{E9EFEF}
%\definecolor{ResultColor}{gray}{0.9}
\DeclareCaptionFont{blue}{\color{captioncolor}}
\captionsetup{labelfont={bf,blue},font={small,blue},textfont={small,blue}}
\providecommand{\lineacross}{\hrule width \textwidth height 0.5pt}
\newcommand{\startexercises}{
 \centerline{\fbox{\textsf{\textbf{\large Exercises}}}}
}
\newcommand{\probs}[1]{\par\noindent\textsf{\textbf{\large #1}}}
\newcommand{\hrulethick}{\leaders\hrule height 5pt depth -2pt\hfill}
\newcommand{\divider}{\noindent{\color{linecolor}\rule{\linewidth}{3pt}}}
\setlength{\multicolsep}{2.5mm}
\setlength{\leftmargini}{1.4em}
\setlength{\shadeboxrule}{.4pt}
\theoremstyle{definition}
\newshadetheorem{defn}{Definition}[chapter]
\newshadetheorem{thm}{Theorem}[chapter]
\newshadetheorem{lem}[thm]{Lemma}
\newshadetheorem{cor}[thm]{Corollary}
\newtheoremstyle{itexmp}%
 {\topsep}
 {\topsep}
 {\normalfont\small}
 {0pt}
 {\itshape\bfseries}{}
 { }
 {#1 \textbf{#2} \enskip {\color{linecolor}\hrulethick}\\[6pt]}
\theoremstyle{itexmp}
\newtheorem{exmp}{Example}[chapter]
\newcommand{\statedefn}[2]{
 \definecolor{shadethmcolor}{HTML}{EDF8FF}
 \definecolor{shaderulecolor}{cmyk}{0.73,0.19,0,0}
 \begin{defn}\label{#1}{#2}\end{defn}
}
\newcommand{\statethm}[2]{
 \definecolor{shadethmcolor}{HTML}{FFFFD8}
 \definecolor{shaderulecolor}{HTML}{9B5900}
 \begin{thm}\label{#1}{#2}\end{thm}
}
\newcommand{\statecor}[2]{
 \definecolor{shadethmcolor}{HTML}{FAF2FF}
 \definecolor{shaderulecolor}{HTML}{9B5900}
 \begin{cor}\label{#1}{#2}\end{cor}
}
\newcommand{\statecomment}[2][0.96\textwidth]{
 \par\noindent\tikzstyle{mybox} = [draw=green!70!black,fill=green!5,thick,rectangle,rounded corners,inner ysep=4pt]
 \begin{tikzpicture}
  \node [mybox] (box){%
   \begin{minipage}{#1}{#2}\end{minipage}
  };
 \end{tikzpicture}
}
\newcommand*\widefbox[1]{\fbox{\hspace{1em}#1\hspace{1em}}}
\renewcommand{\CancelColor}{\color{red}}
\newcommand\myclearpage{\cleartooddpage
 [\thispagestyle{empty}]}
\providecommand{\abs}[1]{\lvert\mspace{1mu}#1\mspace{1mu}\rvert}
\providecommand{\Abs}[1]{\bigl\lvert\mspace{1mu}#1\mspace{1mu}\bigr\rvert}
\providecommand{\norm}[1]{\lVert\mspace{1mu}#1\mspace{1mu}\rVert}
\providecommand{\Norm}[1]{\bigl\lVert\mspace{1mu}#1\mspace{1mu}\bigr\rVert}
\providecommand{\NORM}[1]{\Biggl\lVert\mspace{1mu}#1\mspace{1mu}\Biggr\rVert}
\providecommand{\ssub}[2]{#1_{\scriptscriptstyle #2}}
\providecommand{\ssubsum}[3]{#1_{\scriptscriptstyle #3} + #2_{\scriptscriptstyle #3}}
\providecommand{\Reals}{\mathbb{R}}
\providecommand{\Complex}{\mathbb{C}}
\providecommand{\Rationals}{\mathbb{Q}}
\providecommand{\Naturals}{\mathbb{N}}
\providecommand{\Integers}{\mathbb{Z}}
\providecommand{\Degrees}[0]{\ensuremath{^\circ}}
\providecommand{\ival}[2]{\lbrack #1,#2 \rbrack}

\providecommand{\vectwo}[1]{(#1_{\scriptscriptstyle 1},#1_{\scriptscriptstyle 2})}
\providecommand{\vectwoadd}[2]{(#1_{\scriptscriptstyle 1} + #2_{\scriptscriptstyle 1},
#1_{\scriptscriptstyle 2} + #2_{\scriptscriptstyle 2})}
\providecommand{\vectwosub}[2]{(#1_{\scriptscriptstyle 1} - #2_{\scriptscriptstyle 1},
#1_{\scriptscriptstyle 2} - #2_{\scriptscriptstyle 2})}
\providecommand{\vecthree}[1]{(#1_{\scriptscriptstyle 1},#1_{\scriptscriptstyle 2},#1_{\scriptscriptstyle 3})}
\providecommand{\vecthreeadd}[2]{(#1_{\scriptscriptstyle 1} + #2_{\scriptscriptstyle 1},
#1_{\scriptscriptstyle 2} + #2_{\scriptscriptstyle 2},#1_{\scriptscriptstyle 3} + #2_{\scriptscriptstyle 3})}
\providecommand{\vecthreesub}[2]{(#1_{\scriptscriptstyle 1} - #2_{\scriptscriptstyle 1},
#1_{\scriptscriptstyle 2} - #2_{\scriptscriptstyle 2},#1_{\scriptscriptstyle 3} - #2_{\scriptscriptstyle 3})}
\providecommand{\ssubsum}[3]{#1_{\scriptscriptstyle #3} + #2_{\scriptscriptstyle #3}}
\providecommand{\vecthreeijk}[1]{\ssub{#1}{1}\,\textbf{i} + \ssub{#1}{2}\,\textbf{j} + \ssub{#1}{3}\,\textbf{k}}
\providecommand{\Real}[1]{\mathbb{R}^{#1}}
\providecommand{\Dotprod}[2]{#1 \bm{\cdot} #2}
\providecommand{\Crossprod}[2]{#1 \bm{\times} #2}


\providecommand{\lineintvec}[3]{\int_{#1} \Dotprod{\textbf{#2}}{d\textbf{#3}}}
\providecommand{\olineintvec}[3]{\oint_{#1} \Dotprod{\textbf{#2}}{d\textbf{#3}}}

%\newcommand{\calcbutton}[1]{\setlength\fboxsep{1pt}\cornersize{1}\ovalbox{\small #1}}
%\newcommand{\calcbutton}[1]{\ovalbox{\small #1}}
\newenvironment{proofbar}{%
  \def\FrameCommand{\vrule width 1pt \hspace{5pt}}%
  \MakeFramed {\advance\hsize-\width \FrameRestore}}%
 {\endMakeFramed}
\allowdisplaybreaks[1]
\numberwithin{figure}{section}
\setlength\oddsidemargin{20pt}
\setlength\evensidemargin{20pt}
\setlength\topmargin{0pt}
\setlength\headsep{10pt}
\setlength\textheight{202mm}
%Change indent after chapter number in TOC from 1.5em to 2.2em, and from 2.3 to 3.0 for sections
\makeatletter
\renewcommand*\l@chapter[2]{%
  \ifnum \c@tocdepth >\m@ne
    \addpenalty{-\@highpenalty}%
    \vskip 1.0em \@plus\p@
    \setlength\@tempdima{2.2em}%
    \if@tocleft
      \ifx\toc@l@number\@empty\else
        \setlength\@tempdima{0\toc@l@number}%
      \fi
    \fi
    \begingroup
      \parindent \z@ \rightskip \@pnumwidth
      \parfillskip -\@pnumwidth
      \leavevmode \sectfont
      \advance\leftskip\@tempdima
      \hskip -\leftskip
      #1\nobreak\hfil \nobreak\hb@xt@\@pnumwidth{\hss #2}\par
      \penalty\@highpenalty
    \endgroup
  \fi}
\renewcommand*\l@section{\@dottedtocline{1}{2.2em}{3.0em}}
\makeatother
\makeindex
\makenomenclature



%%%%%%%%%%%%%%%%%%%%%%%%%%%%%%%%%%%%%%%%%%%%%%%%%%%%%%%

%\documentclass[11pt]{book}
\oddsidemargin 0in
\evensidemargin 0in
\marginparwidth 0in
\textheight 8in
\textwidth 6.5in
\topmargin 0in
\usepackage{amssymb,amsmath,amsthm,fancyhdr,supertabular,longtable,hhline}
\usepackage{colortbl}
\usepackage{multirow}  % apparently it was missed by stitz

\usepackage{import,boxedminipage}
\usepackage{chapterfolder}
\usepackage[metapost,truebbox]{mfpic}
%\usepackage[colorlinks, hyperindex, plainpages=false, linkcolor=blue, urlcolor=blue, pdfpagelabels]{hyperref}
\usepackage[all]{hypcap}
\usepackage{cancel}
\usepackage{sectsty}
\usepackage{textcomp}
\allsectionsfont{\mdseries \scshape}
\definecolor{ResultColor}{gray}{0.9}
\theoremstyle{definition}  % this prevents the text in definitions, theorems, and corollaries from being italicized
%\newtheorem{defn}{Definition}[chapter]
%\newtheorem{thm}{Theorem}[chapter]
%\newtheorem{cor}[thm]{Corollary}
\newtheorem{eqn}{Equation}[chapter]
%\newtheorem{ex}{Example}[section]
\setlength{\parindent}{0in}
\newcommand{\bbm}{\begin{boxedminipage}{6.41in}}
\newcommand{\ebm}{\end{boxedminipage}}
\newcounter{HW}
\newcounter{HWindent}
	
\newenvironment{longquote}
	{%
	  	
  		\addtolength{\leftskip}{\egindent}%
  		\small
		\myeqnspacing\protect\sansmath\noindent%
	}
	{%
		 \lmserifmath\par \par\myeqnspacing%
	}


\newenvironment{dialogline}[1]%
	{%
          \noindent{}\sffamily\protect\sansmath%
          {\scshape{#1}: }
        }{%
          \par \par%
          \lmserifmath
	}


%%%%%%%%%%%%%%%%%%%%%%%%%%%%%%%%%%%%%%%%%%%%%%%%%%%%%%%%%
\begin{document}
%%%%%%%%%%%%%%%%%%%%%%%%%%%%%%%%%%%%%%%%%%%%%%%%%%%%%%%%%%%%%%%%%
\deffootnote[1em]{0em}{1em}{\textsuperscript{\thefootnotemark}}
\frontmatter
\pagestyle{empty}
%Put the front cover here
\thisfancypage{\setlength{\fboxsep}{1pt}\setlength\fboxrule{1.5pt}\fbox}{}

\newcommand{\HRule}{\rule{\linewidth}{0.5mm}}
\begin{titlepage}

\begin{center}
\phantom{\rule{5cm}{5.1mm}}
%\phantom{}\\[0.5cm]
% Upper part of the page
   
%\hrule
%{\large \today}

\textsc{\LARGE Physics Program}\\[0.5cm]
\textsc{\LARGE Berea College}\\[0.5cm]
\includegraphics[width=0.20\textwidth]{./preface/BCLogo_BereaBlue_005A8Btr.png}\\[0.5cm]

\HRule \\[0.4cm]
\textsc{\Large PHY-130}\\[0.5cm]

% Title
{ \huge \bfseries Applied Science and Math}\\[0.4cm]

\HRule \\[0.5cm]

% Author and supervisor
\begin{minipage}{0.4\textwidth}
\begin{flushleft} \large
\emph{Year:}\\
Fall \textsc{2020}
\end{flushleft}
\end{minipage}
\begin{minipage}{0.4\textwidth}
\begin{flushright} \large
\emph{Instructor:} \\
Dr.~Martin \textsc{Veillette}
\end{flushright}
\end{minipage}
\vspace{6mm}
\\[1cm]
\input{./preface/oscilloscope.tex}

\vfill

% Bottom of the page
%{\large \today}

\end{center}

\end{titlepage}

%Put the title page here
\title{Applied Mathematics For Engineering and Physics}
\author{\textsf{\textbf{Martin Veillette}}}
\date{\large \textsf{\textsl{Berea College}}}
%Put the author/copyright info here
%\maketitle

\pagestyle{fancy}
\addtokomafont{footnotelabel}{\color{linecolor}}
\addtokomafont{footnotereference}{\color{linecolor}}
\renewcommand{\chaptermark}[1]{\markboth{#1}{}}
\renewcommand{\sectionmark}[1]{\markright{#1}}
\renewcommand{\qedsymbol}{\textsf{\textbf{\textsc{\small{qed}}}}}
\fancyhf{}
\setlength{\headheight}{22.26332pt}
\renewcommand{\headrulewidth}{0pt}
\newlength{\fminilength}%
\setlength{\fminilength}{\textwidth-2\fboxsep-2\fboxrule}%
\fancyhf{}
\fancyhead[LE]{\small\thepage}
\fancyhead[CE]{\small\scshape\leftmark}
\fancyhead[RO]{\small\thepage}
\fancyhead[CO]{\small\scshape\leftmark}
%Put the preface here
%\include{trigbook-preface}
\addchap{Preface}
This book covers the material from the Applied Mathematics in Physics and Engineering course at Berea College. 

This textbook is the result of the effort of many people. It is a compendium of the work that appears in the following open source textbooks: Trigonometry from Micheal Corral, and Light and Matter from Benjamin Crowell.

The Introduction and Scaling chapters are due to Benjamin Crowell from Fullerton College. Micheal Corral from Schoolcraft College contributed to the sections on Vectors in Euclidian Space, General Triangle, Radian Measure,
Right triangle, General Triangles, Identities, Graphing and Inverse functions and Additional Topics.  

Following their work, this book is also released under the GNU Free Documentation License, which allows others to not only copy and distribute the book but also to modify it. Permission is granted to copy, distribute and/or modify this document under the terms of the GNU Free Documentation License, Version 1.2 or any later version published by the Free Software Foundation; with no Invariant Sections, no Front-Cover Texts, and no Back-Cover Texts. A copy of the license is included in the section entitle ``GNU Free Documentation License.''

The original form of this book is LaTeX source code. Compiling this LaTeX source has the effect of generating a device-independent representation of a book, which can be converted to other formats and printed. This book was typeset by the authors using latex, pdflatex, dvips among other free, open-source programs. The graphics were created using MetaPost, TikZ and Gnuplot.

There are numerous exercises throughout the text, which in my experience are more than enough for a semester course in this subject.

Finally, I would like to thank my students in PHY 130 for being the guinea pigs for the initial draft of this book, and for finding the numerous errors and typos it contained.

\begin{flushleft}
\emph{August 2020}\hspace{\stretch{1}}\textsc{Martin Veillette}
\end{flushleft}

%\thispagestyle{empty}
\pdfbookmark[0]{\contentsname}{toc}
%Put the table of contents here
\tableofcontents
%Put a completely blank page here if the TOC has an odd number of pages
\clearpage{\pagestyle{empty}\cleardoublepage}
\mainmatter
\fancyhf{}
\fancyhead[LE]{\tikzstyle{headbox} = [fill=headercolor,line width=0pt,rectangle,rounded corners]
 \begin{tikzpicture}
  \node [headbox] (box){%
   \begin{minipage}{\fminilength}%
     \sffamily{\bfseries\thepage\qquad Chapter \thechapter}\enskip $\bullet$ \enskip\leftmark\hfill{\bfseries\S\thesection}
    \end{minipage}%
  };
 \end{tikzpicture}}
\fancyhead[LO]{\tikzstyle{headbox} = [fill=headercolor,line width=0pt,rectangle,rounded corners]
 \begin{tikzpicture}
  \node [headbox] (box){%
   \begin{minipage}{\fminilength}%
     \sffamily {}\hfill\rightmark\enskip $\bullet$ \enskip{\bfseries{Section \thesection\qquad\thepage}}
    \end{minipage}%
  };
 \end{tikzpicture}}
 
%Put the main chapters here
\chapter{Introduction}
\label{ch:intro}

%\begin{figure}
%\begin{center}
\piccaption{The Mars Climate Orbiter is prepared for its mission.
The laws of physics are the same everywhere, even
on Mars, so the probe could be designed based on
the laws of physics as discovered on earth.
There is unfortunately another reason why this
spacecraft is relevant to the topics of this chapter: it
was destroyed attempting to enter Mars' atmosphere
because engineers at Lockheed Martin forgot to
convert data on engine thrusts from pounds into the
metric unit of force (newtons) before giving the
information to NASA. Conversions are important!}
\parpic[r]{\includegraphics[width=0.5\textwidth]{./introduction/figs/climate-orbiter}
\label{fig:climate-orbiter}}

%\end{center}
%\end{figure}


%%%%%%%%%%%%%%%%%%%%%%%%%%%%%%%%%%%%%%%%%%%%%%%%%%%%%%%%%%%%%%%%
If you drop your shoe and a coin side by side, they hit the
ground at the same time. Why doesn't the shoe get there
first, since gravity is pulling harder on it? How does the
lens of your eye work, and why do your eye's muscles need to
squash its lens into different shapes in order to focus on
objects nearby or far away? These are the kinds of questions
that physics tries to answer about the behavior of light and
matter, the two things that the universe is made of.

\section{The scientific method}

Until very recently in history, no progress was made in
answering questions like these. Worse than that, the
\emph{wrong} answers written by thinkers like the ancient
Greek physicist Aristotle were accepted without question for
thousands of years. Why is it that scientific knowledge has
progressed more since the \index{Renaissance}Renaissance
than it had in all the preceding millennia since the
beginning of recorded history? Undoubtedly the industrial
revolution is part of the answer. Building its centerpiece,
the steam engine, required improved techniques for precise
construction and measurement. (Early on, it was considered a
major advance when English machine shops learned to build
pistons and cylinders that fit together with a gap narrower
than the thickness of a penny.) But even before the
industrial revolution, the pace of discovery had picked up,
mainly because of the introduction of the modern scientific
method. Although it evolved over time, most scientists today
would agree on something like the following list of the
basic principles of the \index{scientific method}scientific method:

%\begin{figure}
%\begin{center}
\piccaption{Science is a cycle of theory and experiment.}
\parpic[r]{\includegraphics{./introduction/figs/thy-and-expt}}

%\end{center}
%\end{figure}

(1) \emph{Science is a cycle of theory and experiment.}
Scientific theories %
% Footnote is to comply with Ca. state standards. Not needed in SN.
\footnote{The term ``theory'' in science does not just mean ``what someone thinks,'' or even ``what a lot of scientists think.'' It means
an interrelated set of statements that have predictive value, and that have survived a broad set of empirical tests. Thus, both Newton's law of gravity and
Darwinian evolution are scientific \emph{theories}. A ``hypothesis,'' in contrast to a theory, is any statement of interest that can be empirically tested.
That the moon is made of cheese is a hypothesis, which was empirically tested, for example, by the Apollo
astronauts.\index{theory}\index{hypothesis}} %
are created to explain the results of
experiments that were created under certain conditions. A
successful theory will also make new predictions about new
experiments under new conditions. Eventually, though, it
always seems to happen that a new experiment comes along,
showing that under certain conditions the theory is not a
good approximation or is not valid at all. The ball is then
back in the theorists' court. If an experiment disagrees
with the current theory, the theory has to be changed, not the experiment.

(2) \emph{Theories should both predict and explain.} The
requirement of predictive power means that a theory is only
meaningful if it predicts something that can be checked
against experimental measurements that the theorist did not
already have at hand. That is, a theory should be testable.
Explanatory value means that many phenomena should be
accounted for with few basic principles. If you answer every
``why'' question with ``because that's the way it is,'' then
your theory has no explanatory value. Collecting lots of
data without being able to find any basic underlying
principles is not science.

(3) \emph{Experiments should be reproducible.} An experiment
should be treated with suspicion if it only works for one
person, or only in one part of the world. Anyone with the
necessary skills and equipment should be able to get the
same results from the same experiment. This implies that
science transcends national and ethnic boundaries; you can
be sure that nobody is doing actual science who claims that
their work is ``Aryan, not Jewish,'' ``Marxist, not
bourgeois,'' or ``Christian, not atheistic.'' An experiment
cannot be reproduced if it is secret, so science is
necessarily a public enterprise.


\begin{figure}[h!]
\begin{center}
  \includegraphics[width=0.5\textwidth]{./introduction/figs/alchemy}
  \caption{A satirical drawing of an alchemist's
      laboratory. H. Cock, after a drawing
      by Peter Brueghel the Elder (16th
      century).}
	  \end{center}
\end{figure}


As an example of the cycle of theory and experiment, a vital
step toward modern chemistry was the experimental observation
that the chemical elements could not be transformed into
each other, e.g., lead could not be turned into gold. This
led to the theory that chemical reactions consisted of
rearrangements of the elements in different combinations,
without any change in the identities of the elements
themselves. The theory worked for hundreds of years, and was
confirmed experimentally over a wide range of pressures and
temperatures and with many combinations of elements. Only in
the twentieth century did we learn that one element could be
trans-formed into one another under the conditions of
extremely high pressure and temperature existing in a
nuclear bomb or inside a star. That observation didn't
completely invalidate the original theory of the immutability
of the elements, but it showed that it was only an
approximation, valid at ordinary temperatures and pressures.


 The scientific method as described here is an idealization,
and should not be understood as a set procedure for doing
science. Scientists have as many weaknesses and character
flaws as any other group, and it is very common for
scientists to try to discredit other people's experiments
when the results run contrary to their own favored point of
view. Successful science also has more to do with luck,
intuition, and creativity than most people realize, and the
restrictions of the scientific method do not stifle
individuality and self-expression any more than the fugue
and sonata forms stifled Bach and Haydn. There is a recent
tendency among social scientists to go even further and to
deny that the scientific method even exists, claiming that
science is no more than an arbitrary social system that
determines what ideas to accept based on an in-group's
criteria. I think that's going too far. If science is an
arbitrary social ritual, it would seem difficult to explain
its effectiveness in building such useful items as
airplanes, CD players, and sewers. If \index{alchemy}alchemy
and \index{astrology}astrology were no less scientific in
their methods than chemistry and astronomy, what was it that
kept them from producing anything useful?



% \startdqswithintro{%
% Consider whether or not the scientific method is being
% applied in the following examples. If the scientific method
% is not being applied, are the people whose actions are being
% described performing a useful human activity, albeit
% an unscientific one?
% }

% \begin{dq}
% Acupuncture is a traditional medical technique of Asian
% origin in which small needles are inserted in the patient's
% body to relieve pain. Many doctors trained in the west
% consider acupuncture unworthy of experimental study because
% if it had therapeutic effects, such effects could not be
% explained by their theories of the nervous system. Who is
% being more scientific, the western or eastern practitioners?
% \end{dq}

% \begin{dq}
% Goethe, a German poet, is less well known for his
% theory of color. He published a book on the subject, in
% which he argued that scientific apparatus for measuring and
% quantifying color, such as prisms, lenses and colored
% filters, could not give us full insight into the ultimate
% meaning of color, for instance the cold feeling evoked by
% blue and green or the heroic sentiments inspired by red. Was
% his work scientific?
% \end{dq}

% \begin{dq}
% A child asks why things fall down, and an adult answers
% ``because of gravity.'' The ancient Greek philosopher
% Aristotle explained that rocks fell because it was their
% nature to seek out their natural place, in contact with the
% earth. Are these explanations scientific?
% \end{dq}

% \begin{dq}
% Buddhism is partly a psychological explanation of human
% suffering, and psychology is of course a science.  The
% Buddha could be said to have engaged in a cycle of theory
% and experiment, since he worked by trial and error, and even
% late in his life he asked his followers to challenge his
% ideas.  Buddhism could also be considered reproducible,
% since the Buddha told his followers they could find
% enlightenment for themselves if they followed a certain
% course of study and discipline.  Is Buddhism a scientific pursuit?
% \end{dq}

\section{What is physics?}


\begin{quotation}
Given for one instant an intelligence which could comprehend
all the forces by which nature is animated and the
respective positions of the things which compose it...nothing
would be uncertain, and the future as the past would be laid
out before its eyes. 
\begin{flushright} - Pierre Simon de Laplace
\end{flushright}
\index{Laplace}
\label{laplace-quote}
\end{quotation}


\index{physics}Physics is the use of the scientific method
to find out the basic principles governing light and matter,
and to discover the implications of those laws. Part of what
distinguishes the modern outlook from the ancient mind-set
is the assumption that there are rules by which the universe
functions, and that those laws can be at least partially
understood by humans. From the Age of Reason through the
nineteenth century, many scientists began to be convinced
that the laws of nature not only could be known but, as
claimed by Laplace, those laws could in principle be used to
predict everything about the universe's future if complete
information was available about the present state of all
light and matter. In subsequent sections, I'll describe two
general types of limitations on prediction using the laws of
physics, which were only recognized in the twentieth century.

\index{matter}Matter can be defined as anything that is
affected by gravity, i.e., that has weight or would have
weight if it was near the Earth or another star or planet
massive enough to produce measurable gravity. \index{light}Light
can be defined as anything that can travel from one place to
another through empty space and can influence matter, but
has no weight. For example, sunlight can influence your body
by heating it or by damaging your DNA and giving you skin
cancer. The physicist's definition of light includes a
variety of phenomena that are not visible to the eye,
including \index{radio waves}radio waves, \index{microwaves}microwaves,
\index{x-rays}x-rays, and \index{gamma rays}gamma rays.
These are the ``colors'' of light that do not happen to fall
within the narrow violet-to-red range of the rainbow that we can see.

Many physical phenomena are not themselves light or matter,
but are properties of light or matter or interactions
between light and matter. For instance, motion is a property
of all light and some matter, but it is not itself light or
matter. The pressure that keeps a bicycle tire blown up is
an interaction between the air and the tire. Pressure is not
a form of matter in and of itself. It is as much a property
of the tire as of the air. Analogously, sisterhood and
employment are relationships among people but are not people themselves.

Some things that appear weightless actually do have weight,
and so qualify as matter. Air has weight, and is thus a form
of matter even though a cubic inch of air weighs less than a
grain of sand. A helium balloon has weight, but is kept from
falling by the force of the surrounding more dense air,
which pushes up on it. Astronauts in orbit around the Earth
have weight, and are falling along a curved arc, but they
are moving so fast that the curved arc of their fall is
broad enough to carry them all the way around the Earth in a
circle. They perceive themselves as being weightless because
their space capsule is falling along with them, and the
floor therefore does not push up on their feet.


The boundary between physics and the other sciences is not
always clear. For instance, chemists study atoms and
molecules, which are what matter is built from, and there
are some scientists who would be equally willing to call
themselves physical chemists or chemical physicists. It
might seem that the distinction between physics and biology
would be clearer, since physics seems to deal with inanimate
objects. In fact, almost all physicists would agree that the
basic laws of physics that apply to molecules in a test tube
work equally well for the combination of molecules that
constitutes a bacterium. (Some might believe that something
more happens in the minds of humans, or even those of cats
and dogs.) What differentiates physics from biology is that
many of the scientific theories that describe living things,
while ultimately resulting from the fundamental laws of
physics, cannot be rigorously derived from physical principles.


\piccaption{Reductionism}
\parpic[r]{\includegraphics{./introduction/figs/reductionism}}

%\begin{figure}
%\begin{center}
%\includegraphics{./introduction/figs/reductionism}
%\caption{Reductionism}
%\end{center}
%\end{figure}


\subsection{Isolated systems and reductionism}

To avoid having to study everything at once, scientists
isolate the things they are trying to study. For instance, a
physicist who wants to study the motion of a rotating
gyroscope would probably prefer that it be isolated from
vibrations and air currents. Even in biology, where field
work is indispensable for understanding how living things
relate to their entire environment, it is interesting to
note the vital historical role played by \index{Darwin}Darwin's
study of the Gal\'{a}pagos Islands, which were conveniently
isolated from the rest of the world. Any part of the
universe that is considered apart from the rest can be
called a ``system.''

\enlargethispage{-2\baselineskip}

Physics has had some of its greatest successes by carrying
this process of isolation to extremes, subdividing the
universe into smaller and smaller parts. Matter can be
divided into atoms, and the behavior of individual atoms can
be studied. Atoms can be split apart into their constituent
neutrons, protons and electrons. Protons and neutrons appear
to be made out of even smaller particles called quarks, and
there have even been some claims of experimental evidence
that quarks have smaller parts inside them. This method of
splitting things into smaller and smaller parts and studying
how those parts influence each other is called \index{reductionism}reductionism.
The hope is that the seemingly complex rules governing the
larger units can be better understood in terms of simpler
rules governing the smaller units. To appreciate what
reductionism has done for science, it is only necessary to
examine a 19th-century chemistry textbook. At that time, the
existence of atoms was still doubted by some, electrons were
not even suspected to exist, and almost nothing was
understood of what basic rules governed the way atoms
interacted with each other in chemical reactions. Students
had to memorize long lists of chemicals and their reactions,
and there was no way to understand any of it systematically.
Today, the student only needs to remember a small set of
rules about how atoms interact, for instance that atoms of
one element cannot be converted into another via chemical
reactions, or that atoms from the right side of the periodic
table tend to form strong bonds with atoms from the left side.

\enlargethispage{-2\baselineskip}

% \startdqs

% \begin{dq}
% I've suggested replacing the ordinary dictionary
% definition of light with a more technical, more precise one
% that involves weightlessness. It's still possible, though,
% that the stuff a lightbulb makes, ordinarily called
% ``light,'' does have some small amount of weight. Suggest an
% experiment to attempt to measure whether it does.
% \end{dq}

% \begin{dq}
% Heat is weightless (i.e., an object becomes no heavier
% when heated), and can travel across an empty room from the
% fireplace to your skin, where it influences you by heating
% you. Should heat therefore be considered a form of light by
% our definition? Why or why not?
% \end{dq}

% \begin{dq}
% Similarly, should sound be considered a form of light?
% \end{dq}

\section{How to learn physics}


\begin{quote}
For as knowledges are now delivered, there is a kind of
contract of error between the deliverer and the receiver;
for he that delivereth knowledge desireth to deliver it in
such a form as may be best believed, and not as may be best
examined; and he that receiveth knowledge desireth rather
present satisfaction than expectant inquiry. - Francis Bacon\index{Bacon, Francis}
\end{quote}

Many students approach a science course with the idea that
they can succeed by memorizing the formulas, so that when a
problem is assigned on the homework or an exam, they will be
able to plug numbers in to the formula and get a numerical
result on their calculator. Wrong! That's not what learning
science is about! There is a big difference between
memorizing formulas and understanding concepts. To start
with, different formulas may apply in different situations.
One equation might represent a definition, which is always
true. Another might be a very specific equation for the
speed of an object sliding down an inclined plane, which
would not be true if the object was a rock drifting down to
the bottom of the ocean. If you don't work to understand
physics on a conceptual level, you won't know which
formulas can be used when.

Most students taking college science courses for the first
time also have very little experience with interpreting the
meaning of an equation. Consider the equation $w=A/h$
relating the width of a rectangle to its height and area. A
student who has not developed skill at interpretation might
view this as yet another equation to memorize and plug in to
when needed. A slightly more savvy student might realize
that it is simply the familiar formula $A=wh$ in a
different form. When asked whether a rectangle would have a
greater or smaller width than another with the same area but
a smaller height, the unsophisticated student might be at a
loss, not having any numbers to plug in on a calculator. The
more experienced student would know how to reason about an
equation involving division --- if $h$ is smaller, and $A$
stays the same, then $w$ must be bigger. Often, students
fail to recognize a sequence of equations as a derivation
leading to a final result, so they think all the intermediate
steps are equally important formulas that they should memorize.

When learning any subject at all, it is important to become
as actively involved as possible, rather than trying to read
through all the information quickly without thinking about
it. It is a good idea to read and think about the questions
posed at the end of each section of these notes as you
encounter them, so that you know you have understood
what you were reading.

Many students' difficulties in physics boil down mainly to
difficulties with math. Suppose you feel confident that you
have enough mathematical preparation to succeed in this
course, but you are having trouble with a few specific
things. In some areas, the brief review given in this
chapter may be sufficient, but in other areas it probably
will not. Once you identify the areas of math in which you
are having problems, get help in those areas. Don't limp
along through the whole course with a vague feeling of dread
about something like scientific notation. The problem will
not go away if you ignore it. The same applies to essential
mathematical skills that you are learning in this course for
the first time, such as vector addition.

Sometimes students tell me they keep trying to understand a
certain topic in the book, and it just doesn't make sense.
The worst thing you can possibly do in that situation is to
keep on staring at the same page. Every textbook explains
certain things badly --- even this one! --- so the best thing to
do in this situation is to look at a different book. Instead
of college textbooks aimed at the same mathematical level as
the course you're taking, you may in some cases find that
high school books or books at a lower math level give
clearer explanations.

 Finally, when reviewing for an exam, don't simply read back
over the text and your lecture notes. Instead, try to use an
active method of reviewing, for instance by discussing some
of the discussion questions with another student, or doing
homework problems you hadn't done the first time.


% end if SN
\section{Self-evaluation}


The introductory part of a book like this is hard to write,
because every student arrives at this starting point with a
different preparation. One student may have grown up
outside the U.S. and so may be completely comfortable with
the metric system, but may have had an algebra course in
which the instructor passed too quickly over scientific
notation. Another student may have already taken  calculus,
but may have never learned the metric system. The following
self-evaluation is a checklist to help you figure out what
you need to study to be prepared for the rest of the course.

\begin{center}
\noindent\begin{tabular}{|p{60mm}|p{60mm}|}
\hline
\textbf{If you disagree with this statement\ldots} & \textbf{you should study this section:}\\
\hline
I am familiar with the basic metric units of meters, kilograms, and seconds, and the most
common metric prefixes: milli- (m), kilo- (k), and centi- (c). & section \ref{sec:metric-basics} Basic of the Metric System \\
\hline
I know about the newton, a unit of force & section \ref{sec:newton-unit} The newton, the Metric Unit of Force\\
\hline
I am familiar with these less common metric prefixes: mega- (M), micro- ($\mu$), and nano- (n). &
section \ref{sec:exotic-prefixes} Less Common Metric Prefixes \\
\hline
I am comfortable with scientific notation. & section \ref{sec:scientific-notation} Scientific Notation \\
\hline
I can confidently do metric conversions. & section \ref{sec:conversions} Conversions \\
\hline
I understand the purpose and use of significant figures. & section \ref{sec:sig-figs} Significant Figures \\
\hline
\end{tabular}
\end{center}
\noindent It wouldn't hurt you to skim the sections you think you
already know about, and to do the self-checks in those sections.

\section{Basics of the metric system}\label{sec:metric-basics}
\index{metric system}

\subsection{The metric system}
\index{metric system}
The range of objects and phenomena studied in physics is immense. From the incredibly short lifetime of a nucleus to the age of the Earth, from the tiny sizes of sub-nuclear particles to the vast distance to the edges of the known universe, from the force exerted by a jumping flea to the force between Earth and the Sun, there are enough factors of 10 to challenge the imagination of even the most experienced scientist. Giving numerical values for physical quantities and equations for physical principles allows us to understand nature much more deeply than does qualitative description alone. To comprehend these vast ranges, we must also have accepted units in which to express them. And we shall find that (even in the potentially mundane discussion of meters, kilograms, and seconds) a profound simplicity of nature appears: all physical quantities can be expressed as combinations of only four fundamental physical quantities: length, mass, time, and electric current. We define a physical quantity either by specifying how it is measured or by stating how it is calculated from other measurements. For example, we define distance and time by specifying methods for measuring them, whereas we define average speed by stating that it is calculated as distance traveled divided by time of travel.
Measurements of physical quantities are expressed in terms of units, which are standardized values. For example, the length of a race, which is a physical quantity, can be expressed in units of meters (for sprinters) or kilometers (for distance runners). Without standardized units, it would be
extremely difficult for scientists to express and compare measured values in a meaningful way.

Units were not standardized until fairly recently in
history, so when the physicist \index{Newton, Isaac}Isaac
Newton gave the result of an experiment with a pendulum, he
had to specify not just that the string was 37  $^7/_8$
inches long but that it was ``37  $^7/_8$ London inches
long.'' The inch as defined in Yorkshire would have been
different. Even after the British Empire standardized its
units, it was still very inconvenient to do calculations
involving money, volume, distance, time, or weight, because
of all the odd conversion factors, like 16 ounces in a
pound, and 5280 feet in a mile. Through the nineteenth
century, schoolchildren squandered most of their mathematical
education in preparing to do calculations such as making
change when a customer in a shop offered a one-crown note
for a book costing two pounds, thirteen shillings and
tuppence. The dollar has always been decimal, and British
money went decimal decades ago, but the United States is
still saddled with the antiquated system of feet, inches,
pounds, ounces and so on.

Every country in the world besides the U.S.~uses a
system of units known in English as the ``metric system.\footnote{Liberia and Myanmar have not legally adopted metric units, but use them
in everyday life.}''
This system is entirely decimal, thanks to the same
eminently logical people who brought about the \index{French
Revolution}French Revolution. In deference to France, the
system's official name is the Syst\`{e}me International, or SI,
meaning International System. (The phrase ``SI system'' is
therefore redundant.)

The wonderful thing about the SI is that people who live in
countries more modern than ours do not need to memorize how
many ounces there are in a pound, how many cups in a pint,
how many feet in a mile, etc. The whole system works with a
single, consistent set of Greek and Latin \index{metric system!prefixes}\index{mega-
(metric prefix)}\index{kilo- (metric prefix)}\index{centi-
(metric prefix)}\index{milli- (metric prefix)}\index{micro-
(metric prefix)}\index{nano- (metric prefix)}prefixes
that modify the basic units. Each
prefix stands for a power of ten, and has an abbreviation
that can be combined with the symbol for the unit. For
instance, the meter is a unit of distance. The prefix kilo-
stands for $10^3$, so a kilometer, 1 km, is a thousand meters.

The basic units of the metric system are the meter for
distance, the second for time, and the gram for mass.

The following are the most common metric prefixes. You
should memorize them.

\begin{center}
\begin{tabular}{llllp{60mm}}
    \multicolumn{2}{c}{prefix} & meaning & \multicolumn{2}{c}{example} \\
    kilo-  &   k   & $10^3$     &   60 kg     & =  a person's mass  \\
    centi- &   c   & $10^{-2}$  &   28 cm     & =  height of a piece of paper  \\
    milli- &   m   & $10^{-3}$  &   1 ms     & =  time for one vibration of a guitar string playing the note D
\end{tabular}
\end{center}

The prefix centi-, meaning $10^{-2}$, is only used in the
centimeter; a hundredth of a gram would not be written as 1
cg but as 10 mg. The centi- prefix can be easily remembered
because a cent is $10^{-2}$  dollars. The official SI
abbreviation for seconds is ``s'' (not ``sec'') and grams
are ``g'' (not ``gm'').

\subsection{The second}


When I stated briefly above that the second was a unit of
time, it may not have occurred to you that this was not
much of a definition. We can make a dictionary-style
definition of a term like ``time,'' or give a general description
like Isaac Newton's: ``Absolute, true, and mathematical time, of itself, and from
its own nature, flows equably without relation to anything external\ldots''
Newton's characterization sounds impressive, but 
physicists today would consider it useless as a definition
of time. Today, the physical sciences are based on
\index{operational definitions}\label{operational-definitions}operational definitions,
which means definitions that spell out the actual steps
(operations) required to measure something numerically.

In an era when our toasters, pens, and coffee pots tell
us the time, it is far from obvious to most people what is
the fundamental operational definition of time. Until
recently, the hour, minute, and second were defined
operationally in terms of the time required for the earth to
rotate about its axis. Unfortunately, the Earth's rotation
is slowing down slightly, and by 1967 this was becoming an
issue in scientific experiments requiring precise time
measurements. The \index{second (unit)}second was therefore
redefined as the time required for a certain number of
vibrations of the light waves emitted by a cesium atoms in a
lamp constructed like a familiar neon sign but with the neon
replaced by cesium. The new definition not only promises to
stay constant indefinitely, but for scientists is a more
convenient way of calibrating a clock than having to carry
out astronomical measurements.




\subsection{The meter}


The French originally defined the meter as $10^{-7}$  times
the distance from the equator to the north pole, as measured
through Paris (of course). Even if the definition was
operational, the operation of traveling to the north pole
and laying a surveying chain behind you was not one that
most working scientists wanted to carry out. Fairly soon, a
standard was created in the form of a metal bar with two
scratches on it. This was replaced by an atomic standard in
1960, and finally in 1983 by the
current definition, which is that
the \index{meter (metric unit)}meter is the
distance traveled by light in a vacuum over a period of
(1/299792458) seconds.

%\begin{figure}
%\begin{center}
%\piccaption[]{The original definition of the meter}
%\parpic[r]{\includegraphics{./introduction/figs/france}}
% \caption{The original definition of the meter}
%\end{center}
%\end{figure}


\subsection{The kilogram}


The third base unit of the SI is the \index{kilogram}kilogram,
a unit of mass.  Mass is intended to be a measure of the
amount of a substance, but that is not an operational
definition. Bathroom scales work by measuring our planet's
gravitational attraction for the object being weighed, but
using that type of scale to define mass operationally would
be undesirable because gravity varies in strength from place
to place on the earth.

There's a surprising amount of disagreement among physics
textbooks about how mass should be defined, but here's how
it's actually handled by the few working physicists who
specialize in ultra-high-precision measurements. They
maintain a physical object in Paris, which is the standard
kilogram, a cylinder made of platinum-iridium alloy.
Duplicates are checked against this mother of all kilograms
by putting the original and the copy on the two opposite
pans of a balance. Although this method of comparison
depends on gravity, the problems associated with differences
in gravity in different geographical locations are bypassed,
because the two objects are being compared in the same
place. The duplicates can then be removed from the Parisian
kilogram shrine and transported elsewhere in the world.
It would be desirable to replace this at some point with a
universally accessible atomic standard rather than one based on
a specific artifact, but as of 2012 the technology for automated counting
of large numbers of atoms has not gotten good
enough to make that work with the desired precision.

\begin{figure}
\begin{center}
\includegraphics[width=0.25\textwidth]{./introduction/figs/standard-kilogram}
\caption{ A duplicate of the Paris kilogram, maintained at the Danish National Metrology Institute.}
\end{center}
\end{figure}





\vfill
\subsection{Combinations of metric units}


Just about anything you want to measure can be measured with
some combination of meters, kilograms, and seconds.  Speed
can be measured in m/s, volume in $\munit^3$, and density in
$\kgunit/\munit^3$. Part of what makes the SI great is this basic
simplicity. No more funny units like a cord of wood, a bolt
of cloth, or a jigger of whiskey. No more liquid and dry
measure. Just a simple, consistent set of units. The SI\index{SI units}
measures put together from meters, kilograms, and seconds
make up the mks system.\index{mks units} For example, the mks unit of
speed is m/s, not km/hr.


\vfill
\subsection{Checking units}


A useful technique for finding mistakes in one's algebra is to analyze
the units associated with the variables.

\pagebreak[4]

\begin{exmp}
\label{eg:checking-units}

Jae starts from the formula $V=\frac{1}{3}Ah$ for the volume of a
cone, where $A$ is the area of its base, and $h$ is its height.
He wants to find an equation that will tell him how tall a conical
tent has to be in order to have a certain volume, given its radius.
His algebra goes like this:
\begin{align}
  V &= \frac{1}{3}Ah \\
  A &= \pi r^2 \\
  V &= \frac{1}{3}\pi r^2 h\\
  h &= \frac{\pi r^2}{3V}
\end{align}
Is his algebra correct? If not, find the mistake.

Line 4 is supposed to be an equation for the height, so the units of the
expression on the right-hand side had better equal meters.
The pi and the 3 are unitless, so we can ignore them.
In terms of units, line 4 becomes
\begin{equation*}
  \munit = \frac{\munit^2}{\munit^3} = \frac{1}{\munit} \qquad .
\end{equation*}
This is false, so there must be a mistake in the algebra. The units
of lines 1, 2, and 3 check out, so the mistake must be in the step
from line 3 to line 4. In fact the result should have been
\begin{equation*}
  h = \frac{3V}{\pi r^2} \qquad .
\end{equation*}
Now the units check: $\munit = \munit^3/\munit^2$.
\end{exmp}
\divider


\begin{center}
\statecomment{\index{Newton, Isaac!definition of time}Isaac Newton wrote,
``\ldots the natural days are truly unequal, though they are
commonly considered as equal, and used for a measure of
time\ldots It may be that there is no such thing as an equable
motion, whereby time may be accurately measured. All motions
may be accelerated or retarded\ldots'' Newton was right. Even
the modern definition of the second in terms of light
emitted by cesium atoms is subject to variation. For
instance, magnetic fields could cause the cesium atoms to
emit light with a slightly different rate of vibration. What
makes us think, though, that a pendulum clock is more
accurate than a sundial, or that a cesium atom is a more
accurate timekeeper than a pendulum clock? That is, how can
one test experimentally how the accuracies of different
time standards compare?
}
\end{center}


\section{The Newton, the metric unit of force}\label{sec:newton-unit}


A force is a push or a pull, or more generally anything that
can change an object's speed or direction of motion. A force
is required to start a car moving, to slow down a baseball
player sliding in to home base, or to make an airplane turn.
(Forces may fail to change an object's motion if they are
canceled by other forces, e.g., the force of gravity pulling
you down right now is being canceled by the force of the
chair pushing up on you.) The metric unit of force is the
Newton, defined as the force which, if applied for one
second, will cause a 1-kilogram object starting from rest to
reach a speed of 1 m/s. 


In section \ref{sec:metric-basics}, I gave a gravitational definition
of mass, but by defining a numerical scale of force, we can
also turn around and define a scale of mass without
reference to gravity. For instance, if a force of two
Newtons is required to accelerate a certain object from rest
to 1 m/s in 1 s, then that object must have a mass of 2
kg. From this point of view, mass characterizes an object's
resistance to a change in its motion, which we call inertia
or inertial mass. Although there is no fundamental reason
why an object's resistance to a change in its motion must be
related to how strongly gravity affects it, careful and
precise experiments have shown that the inertial definition
and the gravitational definition of mass are highly
consistent for a variety of objects. It therefore doesn't
really matter for any practical purpose which definition one adopts.


\section{Less common metric prefixes}\label{sec:exotic-prefixes}


\begin{figure}
\begin{center}
\includegraphics{./introduction/figs/metric-mnemonic}
\caption{      This is a mnemonic to help you remember the most important metric
      prefixes. The word ``little'' is to remind you that the list starts with the prefixes used for
      small quantities and builds upward. The exponent changes by 3, except that of course that we
      do not need a special prefix for $10^0$, which equals one.}
\end{center}
\end{figure}

\vspace{1.5mm}
The following are three metric prefixes which, while less
common than the ones discussed previously, are well worth memorizing.
\begin{center}
\noindent\begin{tabular}{llllp{48mm}}
    \multicolumn{2}{c}{prefix} & meaning & \multicolumn{2}{c}{example} \\
   giga-   &  G      &$10^9$     & 1.4 Gm             &=  diameter of the Sun  \\
    mega-   &  M      &$10^6$     & 6.4 Mm             &=  radius of the Earth  \\
    micro-  &  $\mu$  &$10^{-6}$  & 10 $\mu\munit$     &=  size of a white blood cell  \\
    nano-   &   n     &$10^{-9}$  & 0.154 nm           &=  distance between carbon nuclei in an ethane molecule
\end{tabular}
\end{center}

Note that the abbreviation for micro is the Greek letter mu,
$\mu$ --- a common mistake is to confuse it with m
(milli) or M (mega).

There are other prefixes even less common, used for
extremely large and small quantities.  For instance,
$1\ \text{femtometer}=10^{-15}\ \munit$ is a convenient unit of distance in
nuclear physics, and $1\ \text{gigabyte}=10^9$  bytes is used for
computers' hard disks.  The international committee that
makes decisions about the SI has recently even added some
new prefixes with whimsical names such as  ``yocto-'' $10^{-24}$ and ``zepto-'' $10^{21}$. 

In 2010, Austin Sendek, a physics student at the university of California,  started a petition to establish the prefix ``hella''  as SI prefix for $10^{27}$.  By the time you graduate we may have hellavolts and hellameters !.

\section{Scientific notation}\label{sec:scientific-notation}


Most of the interesting phenomena in our universe
are not on the human scale. It would take about 1,000,000,000,000,000,000,000
bacteria to equal the mass of a human body. When the
physicist Thomas Young discovered that light was a wave, it
was back in the bad old days before scientific notation, and
he was obliged to write that the time required for one
vibration of the wave was 1/500 of a millionth of a
millionth of a second. Scientific notation is a less awkward
way to write very large and very small numbers such as
these.  Here's a quick review.

Scientific notation means writing a number in terms of a
product of something from 1 to 10 and something else that is
a power of ten. For instance,
\begin{align*}
& 32 = 3.2 \times 10^1\\
& 320 =  3.2 \times 10^2\\
& 3200 = 3.2 \times 10^3  \quad\ldots
\end{align*}
Each number is ten times bigger than the previous one.

Since $10^1$  is ten times smaller than $10^2$ , it makes
sense to use the notation $10^0$  to stand for one, the
number that is in turn ten times smaller than $10^1$ .
Continuing on, we can write $10^{-1}$  to stand for 0.1, the
number ten times smaller than $10^0$ . Negative exponents
are used for small numbers:

\begin{align*}
&3.2 =  3.2 \times 10^0\\
&0.32 = 3.2 \times 10^{-1}\\
&0.032 = 3.2 \times 10^{-2} \quad\ldots
\end{align*}

A common source of confusion is the notation used on the
displays of many calculators. Examples:

\vspace{8mm}

\hspace{10mm}\begin{tabular}{ll}
$3.2 \times 10^6$  &     (written notation)\\
3.2E+6             &     (notation on some calculators)\\
$3.2^6$            &     (notation on some other calculators)
\end{tabular}

\vspace{8mm}

\noindent The last example is particularly unfortunate, because
$3.2^6$ really stands for the number 
$3.2 \times 3.2 \times 3.2 \times 3.2 \times 3.2 \times 3.2= 1074$, a totally different number from $3.2 \times 10^6=3200000$.
The calculator notation should never be used in writing.
It's just a way for the manufacturer to save money by
making a simpler display.

\begin{exmp}
A student learns that $10^4$  bacteria, standing in line to
register for classes at Paramecium Community College, would
form a queue of this size:

\includegraphics{./introduction/figs/sc-bacteria-queue-1}

\noindent The student concludes that $10^2$  bacteria would form a line of this length:

\includegraphics{./introduction/figs/sc-bacteria-queue-2}

\noindent Why is the student incorrect?
\end{exmp}
\divider
\vspace{2mm}

\section{Conversions}\label{sec:conversions}
\index{conversions of units}\index{units, conversion of}

I suggest you avoid memorizing lots of conversion factors
between SI units and U.S. units, but two that do come in handy are:

\begin{indentedblock}
 \noindent 1 inch =  2.54 cm
 \noindent An object with a weight on Earth of 2.2 pounds-force has a mass of 1 kg.
\end{indentedblock}

\noindent The first one is the present definition of the inch, so it's exact. The second one is not exact, but is good enough for most purposes. (U.S. units of force and mass are confusing, so it's a good thing they're not used in science. In U.S. units, the unit of force is the pound-force, and the best unit to use for mass is the slug, which is about 14.6 kg.)

More important than memorizing conversion factors is
understanding the right method for doing conversions. Even
within the SI, you may need to convert, say, from grams to
kilograms. Different people have different ways of thinking
about conversions, but the method I'll describe here is
systematic and easy to understand. The idea is that if 1 kg
and 1000 g represent the same mass, then we can consider a fraction like
\begin{equation*}
  \frac{10^3\ \gunit}{1\ \kgunit}
\end{equation*}
to be a way of expressing the number one. This may bother
you. For instance, if you type 1000/1 into your calculator,
you will get 1000, not one. Again, different people have
different ways of thinking about it, but the justification
is that it helps us to do conversions, and it works! Now if
we want to convert 0.7 kg to units of grams, we can multiply
kg by the number one:
\begin{equation*}
  0.7\ \kgunit \times \frac{10^3\ \gunit}{1\ \kgunit}
\end{equation*}
If you're willing to treat symbols such as ``kg'' as if they
were variables as used in algebra (which they're really
not), you can then cancel the kg on top with the kg on the
bottom, resulting in
\begin{equation*}
  0.7\ \cancel{\kgunit} \times \frac{10^3\ \gunit}{1\ \cancel{\kgunit}}  = 700\ \gunit   \qquad   .
\end{equation*}
To convert grams to kilograms, you would simply flip the
fraction upside down.

One advantage of this method is that it can easily be
applied to a series of conversions. For instance, to convert
one year to units of seconds,

\begin{multline*}
1\ \cancel{\text{year}} \times
\frac{365\ \cancel{\text{days}}}{1\ \cancel{\text{year}}} \times
\frac{24\ \cancel{\text{hours}}}{1\ \cancel{\text{day}}} \times
\frac{60\ \cancel{\text{min}}}{1\ \cancel{\text{hour}}} \times
\frac{60\ \sunit}{1\ \cancel{\text{min}}} = 
= 3.15 \times 10^7\ \sunit   \qquad   .
\end{multline*}

\subsection{Should that exponent be positive, or negative?}

A common mistake is to write the conversion fraction
incorrectly. For instance the fraction
\begin{equation*}
  \frac{10^3\ \kgunit}{1\ \gunit} \qquad \text{(incorrect)}
\end{equation*}
does not equal one, because $10^3$  kg is the mass of a car,
and 1 g is the mass of a raisin. One correct way of
setting up the conversion factor would be
\begin{equation*}
  \frac{10^{-3}\ \kgunit}{1\ \gunit} \qquad \text{(correct)} \qquad .
\end{equation*}
You can usually detect such a mistake if you take the time
to check your answer and see if it is reasonable.

If common sense doesn't rule out either a positive or a
negative exponent, here's another way to make sure you get
it right. There are big prefixes and small prefixes:
\begin{center}
\begin{tabular}{ll}
    big prefixes:    & k \quad    M \\
    small prefixes:  &  m  \quad  $\mu$  \quad   n \\
\end{tabular}
\end{center}

\noindent (It's not hard to keep straight which are which, since
``mega'' and ``micro'' are evocative, and it's easy to
remember that a kilometer is bigger than a meter and a
millimeter is smaller.) In the example above, we want the
top of the fraction to be the same as the bottom. Since $k$
is a big prefix, we need to \emph{compensate} by putting a
small number like $10^{-3}$  in front of it, not a big
number like $10^3$.

\begin{exmp}
Each of the following conversions contains an error.  In
each case, explain what the error is.

(a) $1000\ \kgunit \times \frac{1\ \kgunit}{1000\ \gunit}  = 1\ \gunit$

(b) $50\ \munit \times \frac{1\ \zu{cm}}{100\ \munit}    =  0.5\ \zu{cm}$

(c) ``Nano'' is $10^{-9}$, so there are $10^{-9}$  nm in a meter.

(d) ``Micro'' is $10^{-6}$, so 1 kg is $10^6\ \mu\zu{g}$.

\end{exmp}
\divider
\vspace{2mm}


\section{Significant figures}\label{sec:sig-figs}
\index{significant figures}
Science is based on observation and experiment, that is, on measurements. Accuracy is how close a measurement is to the correct value for that measurement. For example, let us say that you are measuring the length of standard computer paper. The packaging in which you purchased the paper states that it is 11.0 inches long. You measure the length of the paper three times and obtain the following measurements: 11.1 in., 11.2 in., and 10.9 in. These measurements are quite accurate because they are very close to the correct value of 11.0 inches. In contrast, if you had obtained a measurement of 12 inches, your measurement would not be very accurate.
The precision of a measurement system is refers to how close the agreement is between repeated measurements (which are repeated under the
same conditions). Consider the example of the paper measurements. The precision of the measurements refers to the spread of the measured
values. One way to analyze the precision of the measurements would be to determine the range, or difference, between the lowest and the highest measured values. In that case, the lowest value was 10.9 in. and the highest value was 11.2 in. Thus, the measured values deviated from each other by at most 0.3 in. These measurements were relatively precise because they did not vary too much in value. However, if the measured values had been 10.9, 11.1, and 11.9, then the measurements would not be very precise because there would be significant variation from one measurement to another.

The degree of accuracy and precision of a measuring system are related to the uncertainty in the measurements. Uncertainty is a quantitative measure of how much your measured values deviate from a standard or expected value. If your measurements are not very accurate or precise, then the uncertainty of your values will be very high. In more general terms, uncertainty can be thought of as a disclaimer for your measured values. For example, if someone asked you to provide the mileage on your car, you might say that it is 45,000 miles, plus or minus 500 miles. The plus or minus amount is the uncertainty in your value. That is, you are indicating that the actual mileage of your car might be as low as 44,500 miles or as high as 45 500 miles, or anywhere in between. All measurements contain some amount of uncertainty. In our example of measuring the length of the paper, we might say that the length of the paper is 11 in., plus or minus 0.2 in. The uncertainty in a measurement, $A$ , is often denoted as $\Delta A$ ("delta A" ),
so the measurement result would be recorded as $A \pm \Delta A$ . In our paper example, the length of the paper could be expressed as $11 \text{in}. \pm 0.2$.

The factors contributing to uncertainty in a measurement include: (1) Limitations of the measuring device, (2) The skill of the person making the measurement, (3) Irregularities in the object being measured, and (4) Any other factors that affect the outcome (highly dependent on the situation). In our example, such factors contributing to the uncertainty could be the following: the smallest division on the ruler is 0.1 in., the person using the ruler has bad eyesight, or one side of the paper is slightly longer than the other. At any rate, the uncertainty in a measurement must be based on a careful consideration of all the factors that might contribute and their possible effects.


Say an engineer is designing a car engine, and has been told
that the diameter of the pistons (which are being designed
by someone else) is 5 cm. He knows that 0.02 cm of clearance
is required for a piston of this size, so he designs the
cylinder to have an inside diameter of 5.04 cm. Luckily, his
supervisor catches his mistake before the car goes into
production. She explains his error to him, and mentally puts
him in the ``do not promote'' category.

What was his mistake? The person who told him the pistons
were 5 cm in diameter was wise to the ways of significant
figures, as was his boss, who explained to him that he
needed to go back and get a more accurate number for the
diameter of the pistons. That person said ``5 cm'' rather
than ``5.00 cm'' specifically to avoid creating the
impression that the number was extremely accurate. In
reality, the pistons' diameter was 5.13 cm. They would never
have fit in the 5.04-cm cylinders.

The number of digits of accuracy in a number is referred to
as the number of significant figures, or ``sig figs'' for
short. As in the example above, sig figs provide a way of
showing the accuracy of a number. In most cases, the result
of a calculation involving several pieces of data can be no
more accurate than the least accurate piece of data. In
other words, ``garbage in, garbage out.'' Since the 5 cm
diameter of the pistons was not very accurate, the result of
the engineer's calculation, 5.04 cm, was really not as
accurate as he thought. In general, your result should not
have more than the number of sig figs in the least accurate
piece of data you started with. The calculation above should
have been done as follows:

\begin{align*}
&5\ \zu{cm} \qquad    \text{(1 sig fig)}\\
+ &0.04\ \zu{cm}  \qquad  \text{(1 sig fig)}\\
= &5\ \zu{cm}   \qquad  \text{(rounded off to 1 sig fig)}
\end{align*}

\noindent The fact that the final result only has one significant
figure then alerts you to the fact that the result is not
very accurate, and would not be appropriate for use in
designing the engine.

Note that the leading zeroes in the number 0.04 do not count
as significant figures, because they are only placeholders.
On the other hand, a number such as 50 cm is ambiguous ---
the zero could be intended as a significant figure, or it
might just be there as a placeholder. The ambiguity
involving trailing zeroes can be avoided by using scientific
notation, in which $5\times10^1$  cm would imply one sig fig of
accuracy, while $5.0\times10^1$  cm would imply two sig figs.

\begin{exmp}
The following quote is taken from an editorial by
Norimitsu Onishi in the New York Times, August 18, 2002.

\noindent 
\begin{quote}
\noindent Consider Nigeria. Everyone agrees it is Africa's most
populous nation. But what is its population? The United
Nations says 114 million; the State Department, 120 million.
The World Bank says 126.9 million, while the Central
Intelligence Agency puts it at 126,635,626.
\end{quote}

\noindent What should bother you about this?
\end{exmp}
\divider

Dealing correctly with significant figures can save you
time! Often, students copy down numbers from their
calculators with eight significant figures of precision,
then type them back in for a later calculation. That's a
waste of time, unless your original data had that kind of
incredible precision.

When combining measurements with different degrees of accuracy and precision, the number of significant digits in the final answer can be no greater than the number of significant digits in the least precise measured value. There are two different rules, one for multiplication and division and the other for addition and subtraction, as discussed below.

1. For multiplication and division: The result should have the same number of significant figures as the quantity having the least significant figures entering into the calculation. For example, the area of a circle can be calculated from its radius using $A = \pi r^2$. Let us see how many significant figures the area has if the radius has only two, say $r = 1.2 m$ . Then,
$$A = \pi r^2 = (3.1415927...)\times (1.2 m)^2 = 4.5238934 m^2.$$
is what you would get using a calculator that has an eight-digit output. But because the radius has only two significant figures, it limits the calculated quantity to two significant figures or
$$A=4.5 m^2,$$ 
even though $\pi$ is good to at least eight digits.

2. For addition and subtraction: The answer can contain no more decimal places than the least precise measurement. Suppose that you buy
7.56-kg of potatoes in a grocery store as measured with a scale with precision 0.01 kg. Then you drop off 6.052-kg of potatoes at your laboratory as measured by a scale with precision 0.001 kg. Finally, you go home and add 13.7 kg of potatoes as measured by a bathroom scale with precision 0.1 kg. How many kilograms of potatoes do you now have, and how many significant figures are appropriate in the answer? The mass is found by simple addition and subtraction: $7.56 kg - 6.052 kg +13.7 kg= 15.208 kg$. 
Next, we identify the least precise measurement: 13.7 kg. This measurement is expressed to the 0.1 decimal place, so our final answer must also be expressed to the 0.1 decimal place. Thus, the answer is rounded to the tenths place, giving us 15.2 kg.

When in doubt, don't use the sig fig rules at all. Instead,
intentionally change one piece of your initial data by the
maximum amount by which you think it could have been off,
and recalculate the final result. The digits on the end that
are completely reshuffled are the ones that are meaningless,
and should be omitted.

\begin{exmp}
How many significant figures are there in each of the
following measurements?

(a). 0.0008

(b). 12,740.0

(c). $5 \times 10^3$

(d). 78.290

(e). 60.62

Solution:

(a) 1; the zeros in this number are placekeepers that indicate the decimal point

(b) 6; here, the zeros indicate that a measurement was made to the 0.1 decimal point, so the zeros are significant

(c) 1; the value $10^3$ signifies the decimal place, not the number of measured values

(d) 5; the final zero indicates that a measurement was made to the 0.001 decimal point, so it is significant

(e) 4; any zeros located in between significant figures in a number are also significant

\end{exmp}
\divider





%%%%%%%%%%%%%%%%%%%%%%%%%%%%%%%%%%%%%%%%%%%%%%%%%%%%%%%%%%%%%%%

% \begin{summary}

% \begin{vocab}

% \vocabitem{matter}{Anything that is affected by gravity.}

% \vocabitem{light}{Anything that can travel from one place to another
% through empty space and can influence matter, but is not
% affected by gravity.}

% \vocabitem{operational definition}{A definition that states what
% operations should be carried out to measure the thing being defined.}

% \vocabitem{Syst\`{e}me International}{A fancy name for the metric system.}

% \vocabitem{mks system}{The use of metric units based on the meter,
% kilogram, and second. Example: meters per second is the mks
% unit of speed, not cm/s or km/hr.}

% \vocabitem{mass}{A numerical measure of how difficult it is to
% change an object's motion.}

% \vocabitem{significant figures}{Digits that contribute to the accuracy of a measurement.}

% \end{vocab}


% \begin{notation}
% \notationitem{m}{meter, the metric distance unit}
% \notationitem{kg}{kilogram, the metric unit of mass}
% \notationitem{s}{second, the metric unit of time}
% \notationitem{M-}{the metric prefix mega-, $10^6$}
% \notationitem{k-}{the metric prefix kilo-, $10^3$}
% \notationitem{m-}{the metric prefix milli-, $10^{-3}$}
% \notationitem{$\mu$-}{the metric prefix micro-, $10^{-6}$}
% \notationitem{n-}{the metric prefix nano-, $10^{-9}$}
% \end{notation}

\section{Summary}

 Physics is the use of the scientific method to study the
 behavior of light and matter. The scientific method requires
 a cycle of theory and experiment, theories with both
 predictive and explanatory value, and reproducible experiments.

 The metric system is a simple, consistent framework for
 measurement built out of the meter, the kilogram, and the
 second plus a set of prefixes denoting powers of ten. The
 most systematic method for doing conversions is shown in
 the following example:
 \begin{equation*}
   370\ \zu{ms} \times \frac{10^{-3}\ \sunit}{1\ \zu{ms}} = 0.37\ \sunit
 \end{equation*}

 Mass is a measure of the amount of a substance. Mass can be 
 defined gravitationally, by comparing an object to a
 standard mass on a double-pan balance, or in terms of
 inertia, by comparing the effect of a force on an object to
 the effect of the same force on a standard mass. The two
 definitions are found experimentally to be proportional to
 each other to a high degree of precision, so we usually
 refer simply to ``mass,'' without bothering to specify which type.

 A force is that which can change the motion of an object.
The metric unit of force is the Newton, defined as the force
 required to accelerate a standard 1-kg mass from rest to a
 speed of 1 m/s in 1 s.

 Scientific notation means, for example, writing $3.2\times10^5$
  rather than 320000.

 Writing numbers with the correct number of significant
 figures correctly communicates how accurate they are. As a
 rule of thumb, the final result of a calculation is no more
 accurate than, and should have no more significant figures
 than, the least accurate piece of data.

% \end{summarytext}

% \end{summary}

%\begin{hwsection}


\divider

\vspace{2mm}

\startexercises\label{sec1dot12}

%%%%% This problem is used by: 1np,0,calculator
%%%%%
\begin{enumerate}
\item  Correct use of a calculator: (a) Calculate $\frac{74658}{53222+97554}$ on a calculator.
[Self-check: The most common mistake results in 97555.40.] 
(b) Which would be more like the price of a TV, and which
would be more like the price of a house, \$$3.5\times10^5$  or \$$3.5^5$?


%%%%% This problem is used by: 1np,0,units
%%%%%
\item Compute the following things. If they don't make sense
because of units, say so.\\
(a) 3 cm + 5 cm 
(b) 1.11 m + 22 cm 
(c) 120 miles + 2.0 hours 
(d) 120 miles / 2.0 hours 


%%%%% This problem is used by: 1np,0,backyard
%%%%%
\item  Your backyard has brick walls on both ends. You measure a
distance of 23.4 m from the inside of one wall to the
inside of the other. Each wall is 29.4 cm thick. How far is
it from the outside of one wall to the outside of the other?
Pay attention to significant figures.


\item 
 The speed of light is $3.0\times10^8$  m/s. Convert
this to furlongs per fortnight.  A furlong is 220 yards, and
a fortnight is 14 days.  An inch is 2.54 cm.

\item
 Express each of the following quantities in micrograms:\\
(a) 10 mg, (b) $10^4$ g, (c) 10 kg, (d) $100 \times 10^3$ g, (e) 1000 ng.


\item Convert 134 mg to units of kg, writing your answer in
scientific notation.

\item 
 In the last century, the average age of the onset of
puberty for girls has decreased by several years. Urban
folklore has it that this is because of hormones fed to beef
cattle, but it is more likely to be because modern girls
have more body fat on the average and possibly because of
estrogen-mimicking chemicals in the environment from the
breakdown of pesticides. A hamburger from a hormone-implanted
steer has about 0.2 ng of estrogen (about double the amount
of natural beef). A serving of peas contains about 300 ng of
estrogen. An adult woman produces about 0.5 mg of estrogen
per day (note the different unit!). (a) How many hamburgers
would a girl have to eat in one day to consume as much
estrogen as an adult woman's daily production? (b) How
many servings of peas? 

\item The usual definition of the mean (average) of two
numbers $a$ and $b$ is $(a+b)/2$. This is called the
arithmetic mean. The geometric mean, however, is defined as
$(ab)^{1/2}$ (i.e., the square root of $ab$). For the sake of definiteness, let's say both
numbers have units of mass. (a) Compute the arithmetic mean
of two numbers that have units of grams. Then convert the
numbers to units of kilograms and recompute their mean. Is
the answer consistent? (b) Do the same for the geometric
mean. (c) If $a$ and $b$ both have units of grams, what
should we call the units of \emph{ab}? Does your answer make
sense when you take the square root? (d) Suppose someone
proposes to you a third kind of mean, called the superduper
mean, defined as $(ab)^{1/3}$. Is this reasonable?

\item 
 In an article on the SARS epidemic, the May 7, 2003 New
York Times discusses conflicting estimates of the disease's
incubation period (the average time that elapses from
infection to the first symptoms). ``The study estimated it
to be 6.4 days. But other statistical calculations ...
showed that the incubation period could be as long as 14.22
days.'' What's wrong here?


\begin{figure}{ht}
\begin{center}
\includegraphics[width=0.25\textwidth]{./introduction/figs/pretzels}
\caption{ Pretzels}
\label{pretzel}
\end{center}
\end{figure}

\vspace{1.5mm}

\item 
 The figure \ref{pretzel} shows the corner of a bag of pretzels. What's wrong here?


\item 
The distance to the horizon is given by the expression $\sqrt{2rh}$, where $r$ is the
radius of the Earth, and $h$ is the observer's height above the Earth's surface. (This can
be proved using the Pythagorean theorem.) Show that the units of this expression make
sense. (See example \ref{eg:checking-units} on p.
\pageref{eg:checking-units} for an example of how to do this.)
Don't try to prove the result, just check its units.

\end{enumerate}

%\end{hwsection}





%\chapter{Scaling}
\label{ch:scaling}

\piccaption{Life would be very different if you were the size of an insect.}{\label{fig:bee-flying}}
\parpic[r]{\includegraphics[width=0.6\textwidth]{./scaling/figs/bee-flying_1.jpg}}

%%%%%%%%%%%%%%%%%%%%%%%%%%%%%%%%%%%%%%%%%%%%%%%%%%%%%%%%%
\section{Introduction}\label{sec:intro-to-scaling}
\vspace{1.5mm}

Why can't an insect be the size of a dog? Some skinny
stretched-out cells in your spinal cord are a meter tall ---
why does nature display no single cells that are not just a
meter tall, but a meter wide, and a meter thick as well?
Believe it or not, these are questions that can be answered
fairly easily without knowing much more about physics than
you already do. The only mathematical technique you really
need is the humble conversion, applied to area and volume.

\piccaption{An unusually large amoeba}
\parpic[r]{\includegraphics[width=0.25\textwidth]{./introduction/figs/big-amoeba}}


\index{area!operational definition}
\subsection{Area and volume}
\index{volume!operational definition}

Area can be defined by saying that we can copy the shape of
interest onto graph paper with 1 cm $\times$ 1 cm squares and
count the number of squares inside. Fractions of squares can
be estimated by eye. We then say the area equals the number
of squares, in units of square cm. Although this might seem
less ``pure'' than computing areas using formulae like
$A=\pi r^2$ for a circle or $A=wh/2$ for a triangle,
those formulae are not useful as definitions of area because
they cannot be applied to irregularly shaped areas.

Units of square cm are more commonly written as $\zu{cm}^2$ in
science. Of course, the unit of measurement symbolized by
``cm'' is not an algebra symbol standing for a number that
can be literally multiplied by itself. But it is advantageous
to write the units of area that way and treat the units as
if they were algebra symbols. For instance, if you have a
rectangle with an area of $6 \munit^2$ and a width of 2 m, then
calculating its length as $(6\ \munit^2)/(2\ \munit)=3\ \munit$ gives a result
that makes sense both numerically and in terms of units.
This algebra-style treatment of the units also ensures that
our methods of converting units work out correctly. For
instance, if we accept the fraction
\begin{equation*}
 \frac{100\ \zu{cm}}{1\ \munit}
\end{equation*}
as a valid way of writing the number one, then one times one
equals one, so we should also say that one can be represented by
\begin{equation*}
 \frac{100\ \zu{cm}}{1\ \munit} \times \frac{100\ \zu{cm}}{1\ \munit} \qquad ,
\end{equation*}
which is the same as
\begin{equation*}
 \frac{10000\ \zu{cm}^2}{1\ \munit^2} \qquad .
\end{equation*}
That means the conversion factor from square meters to
square centimeters is a factor of $10^4$, i.e., a square
meter has $10^4$ square centimeters in it.

All of the above can be easily applied to volume as well,
using one-cubic-centimeter blocks instead of squares on graph paper.

To many people, it seems hard to believe that a square meter
equals 10000 square centimeters, or that a cubic meter
equals a million cubic centimeters --- they think it would
make more sense if there were $100\ \zu{cm}^2$ in $1\ \munit^2$, and 100
$\zu{cm}^3$ in $1\ \munit^3$, but that would be incorrect. The examples
shown in figure \ref{yd2-yd3} aim to make the correct answer
more believable, using the traditional U.S. units of feet
and yards. (One foot is 12 inches, and one yard is three feet.)

\begin{figure}
\centering
\includegraphics[width=0.65\textwidth]{./introduction/figs/yd2-yd3}
\caption{ \qquad Visualizing conversions of area and volume using traditional U.S. units.}
\label{yd2-yd3}
\end{figure}



\begin{exmp}
Based on figure \ref{yd2-yd3}, convince yourself that there are 9
$\zu{ft}^2$ in a square yard, and 27 $\zu{ft}^3$ in a cubic yard,
then demonstrate the same thing symbolically (i.e., with the
method using fractions that equal one).

How many square centimeters are there in a square inch?
(1 inch = 2.54 cm) First find an approximate answer by
making a drawing, then derive the conversion factor more
accurately using the symbolic method.
\end{exmp}
\divider
\vspace{2mm}


% In NP, it's a widecaption, here. SN, the odd/even detection fails, and it's a narrow figure, later on.
\begin{center}
\includegraphics[width=0.4\textwidth]{./introduction/figs/galileo_1.jpg}
\end{center}
%\begin{figure}


\statecomment{Galileo Galilei (1564-1642) was a Renaissance Italian who brought the scientific method to bear on physics, creating the modern version of the science. Coming from a noble but very poor family, Galileo had to drop out of medical school at the University of Pisa when he ran out of money. Eventually becoming a lecturer  in mathematics at the same school, he began a career as a notorious troublemaker by writing a burlesque ridiculing the university's regulations --- he was forced to resign, but found a new teaching position at Padua. He invented the pendulum clock, investigated the motion of falling bodies, and discovered the moons of Jupiter. The thrust of his life's work was to discredit Aristotle's physics by confronting it with contradictory experiments, a program that paved the way for Newton's discovery of the relationship between force and motion.}
%\end{figure}

\section{Scaling of area and volume}\label{sec:scaling}
\index{area!scaling of}\index{scaling}\index{volume!scaling of}

\begin{verse}
Great fleas have lesser fleas\\
Upon their backs to bite 'em.\\
And lesser fleas have lesser still,\\
And so ad infinitum.
\end{verse}\index{Swift, Jonathan}


Now how do these conversions of area and volume relate to
the questions I posed about sizes of living things? Well,
imagine that you are shrunk like Alice in Wonderland to the
size of an insect. One way of thinking about the change of
scale is that what used to look like a centimeter now looks
like perhaps a meter to you, because you're so much smaller.
If area and volume scaled according to most people's
intuitive, incorrect expectations, with $1\ \munit^2$ being the
same as 100 $\zu{cm}^2$, then there would be no particular reason
why nature should behave any differently on your new,
reduced scale. But nature does behave differently now that
you're small. For instance, you will find that you can walk
on water, and jump to many times your own height. The
physicist \index{Galileo Galilei}Galileo Galilei had the
basic insight that the scaling of area and volume determines
how natural phenomena behave differently on different
scales. He first reasoned about mechanical structures, but
later extended his insights to living things, taking the
then-radical point of view that at the fundamental level, a
living organism should follow the same laws of nature as a
machine. We will follow his lead by first discussing
machines and then living things.

\subsection{Galileo on the behavior of nature on large and small scales}

One of the world's most famous pieces of scientific writing
is Galileo's \index{Dialogues Concerning the Two New
Sciences}Dialogues Concerning the Two New Sciences. Galileo
was an entertaining writer who wanted to explain things
clearly to laypeople, and he livened up his work by casting
it in the form of a dialogue among three people. Salviati is
really Galileo's alter ego. Simplicio is the stupid
character, and one of the reasons Galileo got in trouble
with the Church was that there were rumors that Simplicio
represented the \index{Pope}Pope. Sagredo is the earnest and
intelligent student, with whom the reader is supposed to
identify. (The following excerpts are from the 1914
translation by Crew and de Salvio.)

\begin{figure}
\centering
\begin{tabular}{ccc}
\begin{minipage}{150pt}
\frame{\includegraphics[width=0.5\textwidth]{./scaling/figs/boat-1}}
\caption{ The small boat holds up just fine.}
\label{boat-1}
\end{minipage}
&
\begin{minipage}{150pt}
\frame{\includegraphics[width=0.4\textwidth]{./scaling/figs/boat-2}}
\caption{ A larger boat built with the same
 proportions as the small one will
 collapse under its own weight.}
\label{boat-2}
\end{minipage}
&
\begin{minipage}{150pt}
\frame{\includegraphics[width=0.4\textwidth]{./scaling/figs/boat-3}}
\caption{ A boat this large needs to have timbers
 that are thicker compared to its size.}
\label{boat-3}
\end{minipage}
\end{tabular}
\end{figure}

\vspace{1.5mm}


\textbf{Sagredo:} Yes, that is what I mean; and I refer especially to
his last assertion which I have always regarded as false\ldots;
namely, that in speaking of these and other similar machines
one cannot argue from the small to the large, because many
devices which succeed on a small scale do not work on a
large scale. Now, since mechanics has its foundations in
geometry, where mere size [ is unimportant], I do not see
that the properties of circles, triangles, cylinders, cones
and other solid figures will change with their size. If,
therefore, a large machine be constructed in such a way that
its parts bear to one another the same ratio as in a smaller
one, and if the smaller is sufficiently strong for the
purpose for which it is designed, I do not see why the
larger should not be able to withstand any severe and
destructive tests to which it may be subjected.\\


\noindent Salviati contradicts Sagredo:

\textbf{Salviati:} \ldots Please observe, gentlemen, how facts which at
first seem improbable will, even on scant explanation, drop
the cloak which has hidden them and stand forth in naked and
simple beauty. Who does not know that a horse falling from a
height of three or four cubits will break his bones, while a
dog falling from the same height or a cat from a height of
eight or ten cubits will suffer no injury? Equally harmless
would be the fall of a grasshopper from a tower or the fall
of an ant from the distance of the moon.



The point Galileo is making here is that small things are
sturdier in proportion to their size. There are a lot of
objections that could be raised, however. After all, what
does it really mean for something to be ``strong'', to be
``strong in proportion to its size,'' or to be strong ``out
of proportion to its size?'' Galileo hasn't given
operational definitions of things like ``strength,'' i.e.,
definitions that spell out how to measure them numerically.

Also, a cat is shaped differently from a horse --- an
enlarged photograph of a cat would not be mistaken for a
horse, even if the photo-doctoring experts at the National
Inquirer made it look like a person was riding on its back.
A grasshopper is not even a mammal, and it has an exoskeleton
instead of an internal skeleton. The whole argument would be
a lot more convincing if we could do some isolation of
variables, a scientific term that means to change only one
thing at a time, isolating it from the other variables that
might have an effect. If size is the variable whose effect
we're interested in seeing, then we don't really want to
compare things that are different in size but also
different in other ways.


\textbf{Salviati:} \ldots we asked the reason why [shipbuilders] employed
stocks, scaffolding, and bracing of larger dimensions for
launching a big vessel than they do for a small one; and [an
old man] answered that they did this in order to avoid the
danger of the ship parting under its own heavy weight, a
danger to which small boats are not subject?


After this entertaining but not scientifically rigorous
beginning, Galileo starts to do something worthwhile by
modern standards. He simplifies everything by considering
the strength of a wooden plank. The variables involved can
then be narowed down to the type of wood, the width, the
thickness, and the length. He also gives an operational
definition of what it means for the plank to have a certain
strength ``in proportion to its size,'' by introducing the
concept of a plank that is the longest one that would not
snap under its own weight if supported at one end. If you
increased its length by the slightest amount, without
increasing its width or thickness, it would break. He says
that if one plank is the same shape as another but a
different size, appearing like a reduced or enlarged
photograph of the other, then the planks would be strong
``in proportion to their sizes'' if both were just barely
able to support their own weight.


\begin{figure}[h]
\begin{center}
\includegraphics[width=0.8\textwidth]{./scaling/figs/galileo-plank}
\end{center}
\label{galileo-plank}
\caption{ 1. This plank is as long as it can be
 without collapsing under its own weight. If it was a hundredth of an inch
 longer, it would collapse. 2. This plank is made out of the same kind of
 wood. It is twice as thick, twice as long, and twice as wide. It will
 collapse under its own weight.}
\end{figure}

Also, Galileo is doing something that would be frowned on in
modern science: he is mixing experiments whose results he
has actually observed (building boats of different sizes),
with experiments that he could not possibly have done
(dropping an ant from the height of the moon).
He now relates how he has done actual experiments with such
planks, and found that, according to this operational
definition, they are not strong in proportion to their
sizes. The larger one breaks. He makes sure to tell the
reader how important the result is, via Sagredo's astonished response:

\textbf{Sagredo:} My brain already reels. My mind, like a cloud
momentarily illuminated by a lightning flash, is for an
instant filled with an unusual light, which now beckons to
me and which now suddenly mingles and obscures strange,
crude ideas. From what you have said it appears to me
impossible to build two similar structures of the same
material, but of different sizes and have them proportionately strong.

In other words, this specific experiment, using things like
wooden planks that have no intrinsic scientific interest,
has very wide implications because it points out a general
principle, that nature acts differently on different scales.


To finish the discussion, Galileo gives an explanation. He
says that the strength of a plank (defined as, say, the
weight of the heaviest boulder you could put on the end
without breaking it) is proportional to its cross-sectional
area, that is, the surface area of the fresh wood that would
be exposed if you sawed through it in the middle. Its
weight, however, is proportional to its volume.
\footnote{Galileo
makes a slightly more complicated argument, taking into account
the effect of leverage (torque). The result I'm
referring to comes out the same regardless of this effect.}

How do the volume and cross-sectional area of the longer
plank compare with those of the shorter plank? We have
already seen, while discussing conversions of the units of
area and volume, that these quantities don't act the way
most people naively expect. You might think that the volume
and area of the longer plank would both be doubled compared
to the shorter plank, so they would increase in proportion
to each other, and the longer plank would be equally able to
support its weight. You would be wrong, but Galileo knows
that this is a common misconception, so he has Salviati
address the point specifically:

\textbf{Salviati:} \ldots Take, for example, a cube two inches on a side
so that each face has an area of four square inches and the
total area, i.e., the sum of the six faces, amounts to
twenty-four square inches; now imagine this cube to be sawed
through three times [with cuts in three perpendicular
planes] so as to divide it into eight smaller cubes, each
one inch on the side, each face one inch square, and the
total surface of each cube six square inches instead of
twenty-four in the case of the larger cube. It is evident
therefore, that the surface of the little cube is only
one-fourth that of the larger, namely, the ratio of six to
twenty-four; but the volume of the solid cube itself is only
one-eighth; the volume, and hence also the weight,
diminishes therefore much more rapidly than the surface\ldots
You see, therefore, Simplicio, that I was not mistaken when
\ldots I said that the surface of a small solid is comparatively
greater than that of a large one.

The same reasoning applies to the planks. Even though they
are not cubes, the large one could be sawed into eight small
ones, each with half the length, half the thickness, and
half the width. The small plank, therefore, has more surface
area in proportion to its weight, and is therefore able to
support its own weight while the large one breaks.

\piccaption{ Galileo discusses planks made of
 wood, but the concept may be easier
 to imagine with clay. }
\parpic[r]{\includegraphics[width=0.25\textwidth]{./scaling/figs/clay} \label{clay}}


All three clay rods
 in the figure \ref{clay} were originally the same
 shape. The medium-sized one was
 twice the height, twice the length, and
 twice the width of the small one, and
 similarly the large one was twice as
 big as the medium one in all its linear
 dimensions. The big one has four
 times the linear dimensions of the
 small one, 16 times the cross-sectional
 area when cut perpendicular to the
 page, and 64 times the volume. That
 means that the big one has 64 times
 the weight to support, but only 16 times
 the strength compared to the smallest
 one.



\vspace{1.5mm}

\subsection{Scaling of area and volume for irregularly shaped objects}


You probably are not going to believe Galileo's claim that
this has deep implications for all of nature unless you can
be convinced that the same is true for any shape. Every
drawing you've seen so far has been of squares, rectangles,
and rectangular solids. Clearly the reasoning about sawing
things up into smaller pieces would not prove anything
about, say, an egg, which cannot be cut up into eight
smaller egg-shaped objects with half the length.


\piccaption{\quad Violin }
\parpic[l]{\includegraphics[width=0.2\textwidth]{./scaling/figs/violin} \label{violin}}



\vspace{1.5mm}

Is it always true that something half the size has one
quarter the surface area and one eighth the volume, even if
it has an irregular shape? Take the example of a child's
violin. Violins are made for small children in smaller size
to accomodate their small bodies. Figure \ref{violin} shows
a full-size violin, along with two
violins made with half and 3/4 of the normal length.\footnote{The customary
terms ``half-size'' and ``3/4-size'' actually don't describe the
sizes in any accurate way. They're really just standard, arbitrary
marketing labels.}
Let's study the surface area of the front
panels of the three violins.

Consider the square in the interior of the panel of the
full-size violin. In the 3/4-size violin, its height and
width are both smaller by a factor of 3/4, so the area of
the corresponding, smaller square becomes $3/4\times3/4=9/16$ of
the original area, not 3/4 of the original area. Similarly,
the corresponding square on the smallest violin has half the
height and half the width of the original one, so its area
is 1/4 the original area, not half.

The same reasoning works for parts of the panel near the
edge, such as the part that only partially fills in the
other square. The entire square scales down the same as a
square in the interior, and in each violin the same fraction
(about 70\%) of the square is full, so the contribution of
this part to the total area scales down just the same.


Since any small square region or any small region covering
part of a square scales down like a square object, the
entire surface area of an irregularly shaped object changes
in the same manner as the surface area of a square: scaling
it down by 3/4 reduces the area by a factor of 9/16, and so on.

In general, we can see that any time there are two objects
with the same shape, but different linear dimensions (i.e.,
one looks like a reduced photo of the other), the ratio of
their areas equals the ratio of the squares of their linear dimensions:
\begin{equation*}
 \frac{A_1}{A_2} = \left(\frac{L_1}{L_2}\right)^2 \qquad .
\end{equation*}
Note that it doesn't matter where we choose to measure the
linear size, $L$, of an object. In the case of the violins,
for instance, it could have been measured vertically,
horizontally, diagonally, or even from the bottom of the
left f-hole to the middle of the right f-hole. We just have
to measure it in a consistent way on each violin. Since all
the parts are assumed to shrink or expand in the same
manner, the ratio $L_1/L_2$ is independent of the choice of measurement.


\piccaption{
 The muffin comes out of the oven too hot to eat. Breaking
 it up into four pieces increases its surface area while keeping the total
 volume the same. It cools faster because of the greater surface-to-volume ratio.}
\parpic[r]{\includegraphics[width=0.4\textwidth]{./scaling/figs/muffin} \label{muffin}}
 


\vspace{1.5mm}

It is also important to realize that it is completely
unnecessary to have a formula for the area of a violin. It
is only possible to derive simple formulas for the areas of
certain shapes like circles, rectangles, triangles and so
on, but that is no impediment to the type of reasoning we are using.

Sometimes it is inconvenient to write all the equations in
terms of ratios, especially when more than two objects are
being compared. A more compact way of rewriting the
previous equation is
\begin{equation*}
 A \propto L^2 \qquad .
\end{equation*}
The symbol ``$\propto$'' means ``is proportional to.''
Scientists and engineers often speak about such relationships
verbally using the phrases ``scales like'' or ``goes like,''
for instance ``area goes like length squared.''

All of the above reasoning works just as well in the case of
volume. Volume goes like length cubed:
\begin{equation*}
 V \propto L^3 \qquad .
\end{equation*}

\begin{selfcheck}{axle-ratios}
When a car or truck travels over a road, there is wear and tear on the road surface, which incurs a cost.
Studies show that the cost $C$ per kilometer of travel is related to the weight per axle $w$
by $C \propto w^4$.
Translate this into a statement about ratios.
\end{selfcheck}



If different objects are made of the same material with the
same density, $\rho =m/V$, then their masses, $m=\rho V$,
are proportional to $L^3$. (The
symbol for density is $\rho$, the lower-case Greek letter ``rho.'')

An important point is that all of the above reasoning about
scaling only applies to objects that are the same shape. For
instance, a piece of paper is larger than a pencil, but has
a much greater surface-to-volume ratio.

\begin{exmp}
\label{eg:scale-triangle} 
In figure \ref{eg-scale-triangle-1}, the larger triangle has sides twice as long. How many times greater is its area?

Correct solution \#1: Area scales in proportion to the
square of the linear dimensions, so the larger triangle has
four times more area $(2^2=4)$.

\begin{figure}[h]
\label{eg-scale-triangle-1} 
\begin{center}
\begin{tikzpicture}
    \path[coordinate] (0,0)  coordinate(A)
                ++( 90:2cm) coordinate(B)
                ++(-45:2.8cm) coordinate(C);
    \draw[fill=blue!60,opacity=0.3] (A) -- (B) -- (C) -- cycle;

 \path[coordinate] (5,0)  coordinate(D)
                ++( 90:4cm) coordinate(E)
                ++(-45:5.6cm) coordinate(F);
    \draw[fill=blue!60,opacity=0.3] (D) -- (E) -- (F) -- cycle;
    
\end{tikzpicture}
\caption{ \qquad
 Example \ref{eg:scale-triangle}. The big triangle has four times more
 area than the little one.}
\end{center}
\end{figure}

\begin{figure}[h]
\label{eg-scale-triangle-1} 
\begin{center}
\begin{tikzpicture}
    \def\couleur{blue}
    \path[coordinate] (0,0)  coordinate(A)
                ++( 90:4cm) coordinate(B)
                ++(-45:5.6cm) coordinate(C);
 \draw[fill=blue!60,opacity=0.3] (A) -- (B) -- (C) -- cycle;

 \path[coordinate] (0,2)  coordinate(D)
                ++( 0:2cm) coordinate(E)
                ++(-90:2cm) coordinate(F);
    \draw[thick,dash pattern=on5pt off3pt] (D) -- (E) -- (F) -- cycle;
    
\end{tikzpicture}
\caption{ \qquad
 A tricky way of solving example \ref{eg:scale-triangle}, explained in
 solution \#2.}
\end{center}
\end{figure}



\vspace{1.5mm}

\end{exmp}

Correct solution \#2: You could cut the larger triangle into
four of the smaller size, as shown in fig. (b), so its area
is four times greater. (This solution is correct, but it
would not work for a shape like a circle, which can't be cut
up into smaller circles.)

Correct solution \#3: The area of a triangle is given by

$A=bh/2$, where $b$ is the base and $h$ is the height. The
areas of the triangles are
\begin{align*}
 A_1 &= b_1 h_1/2\\
 A_2 &= b_2 h_2/2\\
 &= (2b_1)(2h_1)/2 \\
 &= 2b_1 h_1\\
 A_2/A_1 &=(2b_1 h_1)/(b_1 h_1/2)\\
 &= 4
\end{align*}

(Although this solution is correct, it is a lot more work
than solution \#1, and it can only be used in this case
because a triangle is a simple geometric shape, and we
happen to know a formula for its area.)

Correct solution \#4: The area of a triangle is $A= bh/2$.
The comparison of the areas will come out the same as long
as the ratios of the linear sizes of the triangles is as
specified, so let's just say $b_1=1.00$ m and $b_2=2.00$
m. The heights are then also $h_1=1.00$ m and $h_2=2.00$
m, giving areas $A_1=0.50\ \munit^2$ and $A_2=2.00\ \munit^2$, so $A_2/A_1=4.00$.

(The solution is correct, but it wouldn't work with a shape
for whose area we don't have a formula. Also, the numerical
calculation might make the answer of 4.00 appear inexact,
whereas solution \#1 makes it clear that it is exactly 4.)

\enlargethispage{\baselineskip}

Incorrect solution: The area of a triangle is $A=bh/2$, and
if you plug in $b=2.00$ m and $h=2.00$ m, you get $A=2.00\ \munit^2$,
so the bigger triangle has 2.00 times more area. (This
solution is incorrect because no comparison has been made
with the smaller triangle.)

\vfill\pagebreak[4]

\begin{exmp}\label{eg:scale-sphere}
In figure \ref{eg-scale-sphere}, the larger sphere has a radius that
is five times greater. How many times greater is its volume?




\usetikzlibrary{calc,fadings,decorations.pathreplacing}



%% document-wide tikz options and styles

\tikzset{%
  >=latex, % option for nice arrows
  inner sep=0pt,%
  outer sep=2pt,%
  mark coordinate/.style={inner sep=0pt,outer sep=0pt,minimum size=3pt,
    fill=black,circle}%
}


\begin{figure}
\begin{center}
\begin{tikzpicture} % "THE GLOBE" showcase

\def\RL{2.5} % sphere radius
\def\RS{0.5} % sphere radius
\def\angEl{35} % elevation angle
\filldraw[ball color=white] (5,0) circle (\RL);
\filldraw[ball color=white] (0,-2) circle (\RS);
%\foreach \t in {-80,-60,...,80} { \DrawLatitudeCircle[\R]{\t} }
%\foreach \t in {-5,-35,...,-175} { \DrawLongitudeCircle[\R]{\t} }
\end{tikzpicture}
\end{center}
\caption{ \qquad  \qquad  Example \ref{eg:scale-sphere}: 
 The big sphere has 125 times more
 volume than the little one.
 }
\end{figure}


\vspace{1.5mm}

Correct solution \#1: Volume scales like the third power of
the linear size, so the larger sphere has a volume that is
125 times greater $(5^3=125)$.

Correct solution \#2: The volume of a sphere is $V=(4/3)\pi r^3$, so
\begin{align*}
 V_1 &= \frac{4}{3}\pi r_1^3 \\
 V_2 &= \frac{4}{3}\pi r_2^3 \\
 &= \frac{4}{3}\pi (5r_1)^3 \\
 &= \frac{500}{3}\pi r_1^3 \\
 V_2/V_1 &= \left( \frac{500}{3}\pi r_1^3 \right) / \left( \frac{4}{3}\pi r_1^3 \right)
 &= 125
\end{align*}

Incorrect solution: The volume of a sphere is $V=(4/3)\pi r^3$, so
\begin{align*}
 V_1 &= \frac{4}{3}\pi r_1^3 \\
 V_2 &= \frac{4}{3}\pi r_2^3 \\
 &= \frac{4}{3}\pi \cdot 5r_1^3 \\
 &= \frac{20}{3}\pi r_1^3 \\
 V_2/V_1 &= \left( \frac{20}{3}\pi r_1^3 \right) / \left( \frac{4}{3}\pi r_1^3 \right)
 &= 5
\end{align*}

(The solution is incorrect because $(5r_1)^3$ is not the same as $5r_1^3$.)


\end{exmp}

\begin{figure}[h] 
\begin{center}
\begin{tikzpicture}
\draw[color=blue] (0,0) node[below] {\fontsize{48}{48}\selectfont S};
\draw[color=blue] (2,-0.25) node[below] {\fontsize{36}{36}\selectfont S};
\end{tikzpicture}

%\includegraphics[width=0.2\textwidth]{./scaling/figs/eg-scale-letter-s}
\caption{ \qquad %
%
 Example \ref{eg:scale-letter-s}. The 48-point ``S'' has 1.78 times
 more area than the 36-point ``S.''
 }
\end{center}
\label{eg-scale-letter-s}
\end{figure}


\vspace{1.5mm}

\begin{exmp}
\label{eg:scale-letter-s}
 The first letter ``S'' in figure \ref{eg-scale-letter-s} is in a
36-point font, the second in 48-point. How many times more
ink is required to make the larger ``S''? (Points are a unit
of length used in typography.)

Correct solution: The amount of ink depends on the area to
be covered with ink, and area is proportional to the square
of the linear dimensions, so the amount of ink required for
the second ``S'' is greater by a factor of $(48/36)^2=1.78$.

Incorrect solution: The length of the curve of the second
``S'' is longer by a factor of $48/36=1.33$, so 1.33 times
more ink is required.

(The solution is wrong because it assumes incorrectly that
the width of the curve is the same in both cases. Actually
both the width and the length of the curve are greater by a
factor of 48/36, so the area is greater by a factor
of $(48/36)^2=1.78$.)
\end{exmp}
\divider
\vspace{2mm}

Reasoning about ratios and proportionalities
is one of the three essential mathematical skills, summarized on pp.\pageref{begin-skills}-\pageref{end-skills},
that you need for success in this course.


Galileo spends a lot of time in his dialog discussing
what really happens when things break. He discusses
everything in terms of Aristotle's now-discredited
explanation that things are hard to break, because if
something breaks, there has to be a gap between the two
halves with nothing in between, at least initially. Nature,
according to Aristotle, ``abhors a vacuum,'' i.e., nature
doesn't ``like'' empty space to exist. Of course, air will
rush into the gap immediately, but at the very moment of
breaking, Aristotle imagined a vacuum in the gap. Is
Aristotle's explanation of why it is hard to break things an
experimentally testable statement? If so, how could it be
tested experimentally?
\





\section{Scaling applied to biology}
\index{scaling!applied to biology}

\subsection{Organisms of different sizes with the same shape}


The left-hand panel in figure \ref{scaling-animals-1} 
shows the approximate
validity of the proportionality $m\propto L^3$ for \index{cockroaches}cockroaches
(redrawn from McMahon and Bonner). The scatter of the
points around the curve indicates that some cockroaches are
proportioned slightly differently from others, but in
general the data seem well described by $m\propto L^3$. That
means that the largest cockroaches the experimenter could
raise (is there a 4-H prize?) had roughly the same shape
as the smallest ones.

\begin{figure}
\begin{center}
\includegraphics[width=0.7\textwidth]{./scaling/figs/scaling-animals-1}
\caption{ \qquad Geometrical scaling of animals.}
\label{scaling-animals-1}
\end{center}
\end{figure}


Another relationship that should exist for animals of
different sizes shaped in the same way is that between
surface area and body mass. If all the animals have the same
average density, then body mass should be proportional to
the cube of the animal's linear size, $m\propto L^3$, while
surface area should vary proportionately to $L^2$.
Therefore, the animals' surface areas should be proportional
to $m^{2/3}$. As shown in the right-hand panel of figure
\ref{scaling-animals-1}, this relationship
appears to hold quite well for the dwarf siren, a type of
\index{salamanders}salamander. Notice how the curve bends
over, meaning that the surface area does not increase as
quickly as body mass, e.g., a salamander with eight times
more body mass will have only four times more surface area.

\begin{figure}
\begin{center}
\includegraphics[width=0.7\textwidth]{./scaling/figs/scaling-animals-2}
\caption{ \qquad Scaling of animals' bodies related to metabolic rate and skeletal strength.}
\end{center}
\end{figure}


This behavior of the ratio of surface area to mass (or,
equivalently, the ratio of surface area to volume) has
important consequences for mammals, which must maintain a
constant body temperature. It would make sense for the rate
of heat loss through the animal's skin to be proportional to
its surface area, so we should expect small animals, having
large ratios of surface area to volume, to need to produce a
great deal of heat in comparison to their size to avoid
dying from low body temperature. This expectation is borne
out by the data of the left-hand panel of
figure \ref{scaling-animals-2}, showing the rate of
oxygen consumption of guinea pigs as a function of their
body mass. Neither an animal's heat production nor its
surface area is convenient to measure, but in order to
produce heat, the animal must metabolize oxygen, so oxygen
consumption is a good indicator of the rate of heat
production. Since surface area is proportional to $m^{2/3}$, the
proportionality of the rate of oxygen consumption to $m^{2/3}$ is
consistent with the idea that the animal needs to produce
heat at a rate in proportion to its surface area. Although
the smaller animals metabolize less oxygen and produce less
heat in absolute terms, the amount of food and oxygen they
must consume is greater in proportion to their own mass. The
Etruscan pigmy shrew, weighing in at 2 grams as an adult, is
at about the lower size limit for mammals. It must eat
continually, consuming many times its body weight each day to survive.

\subsection{Changes in shape to accommodate changes in size}


Large mammals, such as \index{elephant}elephants, have a
small ratio of surface area to volume, and have problems
getting rid of their heat fast enough. An elephant cannot
simply eat small enough amounts to keep from producing
excessive heat, because cells need to have a certain minimum
metabolic rate to run their internal machinery. Hence the
elephant's large ears, which add to its surface area and
help it to cool itself. Previously, we have seen several
examples of data within a given species that were consistent
with a fixed shape, scaled up and down in the cases of
individual specimens. The elephant's ears are an example of
a change in shape necessitated by a change in scale.

\begin{figure}[h] 
\begin{center}
\includegraphics[width=0.4\textwidth]{./scaling/figs/galileo-bones}
\caption{ \qquad  Galileo's original drawing, showing
 how larger animals' bones must be
 greater in diameter compared to their
 lengths.
 }
\end{center}
\end{figure}


\vspace{1.5mm}
Large animals also must be able to support their own weight.
Returning to the example of the strengths of planks of
different sizes, we can see that if the strength of the
plank depends on area while its weight depends on volume,
then the ratio of strength to weight goes as follows:
\begin{equation*}
 \text{strength}/\text{weight} \propto A/V \propto 1/L \qquad .
\end{equation*}
Thus, the ability of objects to support their own weights
decreases inversely in proportion to their linear dimensions.
If an object is to be just barely able to support its own
weight, then a larger version will have to be proportioned
differently, with a different shape.

Since the data on the cockroaches seemed to be consistent
with roughly similar shapes within the species, it appears
that the ability to support its own weight was not the
tightest design constraint that Nature was working under
when she designed them. For large animals, structural
strength is important. Galileo was the first to quantify
this reasoning and to explain why, for instance, a large
animal must have bones that are thicker in proportion to
their length. Consider a roughly cylindrical bone such as a
leg bone or a \index{vertebra}vertebra. The length of the
bone, $L$, is dictated by the overall linear size of the
animal, since the animal's skeleton must reach the animal's
whole length. We expect the animal's mass to scale as $L^3$,
so the strength of the bone must also scale as $L^3$.
Strength is proportional to cross-sectional area, as with
the wooden planks, so if the diameter of the bone is $d$, then
\begin{align*}
 d^2 &\propto L^3 \\
\intertext{or}
 d &\propto L^{3/2} \qquad .
\end{align*}
If the shape stayed the same regardless of size, then all
linear dimensions, including $d$ and $L$, would be
proportional to one another. If our reasoning holds, then
the fact that $d$ is proportional to $L^{3/2}$, not $L$, implies
a change in proportions of the bone. As shown in the right-hand
panel of figure \ref{scaling-animals-2},
the vertebrae of African Bovidae
follow the rule $d\propto L^{3/2}$ fairly well. The vertebrae
of the giant eland are as chunky as a coffee mug, while
those of a Gunther's dik-dik are as slender as the cap of a pen.



\statecomment{
Single-celled animals must passively absorb nutrients and
oxygen from their surroundings, unlike humans who have lungs
to pump air in and out and a heart to distribute the
oxygenated blood throughout their bodies. Even the cells
composing the bodies of multicellular animals must absorb
oxygen from a nearby capillary through their surfaces. Based
on these facts, explain why cells are always microscopic in size.


The reasoning of the previous question would seem to be
contradicted by the fact that human nerve cells in the
spinal cord can be as much as a meter long, although their
widths are still very small. Why is this possible?
}


% end conditional, biology application only for LM, not SN


\section{Order-of-magnitude estimates}
\index{order-of-magnitude estimates}

\begin{quotation}
It is the mark of an instructed mind to rest satisfied with
the degree of precision that the nature of the subject
permits and not to seek an exactness where only an
approximation of the truth is possible. ~Aristotle
\end{quotation}

It is a common misconception that science must be exact. For
instance, in the Star Trek TV series, it would often happen
that Captain Kirk would ask Mr. Spock, ``Spock, we're in a
pretty bad situation. What do you think are our chances of
getting out of here?'' The scientific Mr. Spock would answer
with something like, ``Captain, I estimate the odds as
237.345 to one.'' In reality, he could not have estimated
the odds with six significant figures of accuracy, but
nevertheless one of the hallmarks of a person with a good
education in science is the ability to make estimates that
are likely to be at least somewhere in the right ballpark.

On many occasions, physicists, other scientists, and engineers need to make approximations or “guesstimates” for a particular quantity. What is the distance to a certain destination? What is the approximate density of a given item? About how large a current will there be in a circuit? Many approximate numbers are based on formulae in which the input quantities are known only to a limited accuracy. As you develop problem-solving skills (that can be applied to a variety of fields through a study of physics), you will also develop skills at approximating. You will develop these skills through thinking more quantitatively, and by being willing to take risks. As with any endeavor, experience helps, as well as familiarity with units. These approximations allow us to rule out certain scenarios or unrealistic numbers. Approximations also allow us to challenge others and guide us in our approaches to our scientific world. Let us do two examples to illustrate this concept.

In many such situations, it is often only necessary to get
an answer that is off by no more than a factor of ten in
either direction. Since things that differ by a factor of
ten are said to differ by one order of magnitude, such an
estimate is called an order-of-magnitude estimate. The
tilde, $\sim$, is used to indicate that things are only of
the same order of magnitude, but not exactly equal, as in
\begin{equation*}
 \text{odds of survival} \sim \text{100 to one} \qquad .
\end{equation*}
The tilde can also be used in front of an individual number
to emphasize that the number is only of the right order of magnitude.

Although making order-of-magnitude estimates seems simple
and natural to experienced scientists, it's a mode of
reasoning that is completely unfamiliar to most college
students. Some of the typical mental steps can be illustrated
in the following example.

\begin{exmp}
Cost of transporting tomatoes (incorrect solution):
Roughly what percentage of the price of a tomato
comes from the cost of transporting it in a truck?

The following incorrect solution illustrates one of the main
ways you can go wrong in order-of-magnitude estimates.

Incorrect solution: Let's say the trucker needs to make a
\$400 profit on the trip. Taking into account her benefits,
the cost of gas, and maintenance and payments on the truck,
let's say the total cost is more like \$2000. I'd guess
about 5000 tomatoes would fit in the back of the truck, so
the extra cost per tomato is 40 cents. That means the cost
of transporting one tomato is comparable to the cost of the
tomato itself. Transportation really adds a lot to the cost
of produce, I guess.
\end{exmp}
\divider
\vspace{2mm}

The problem is that the human brain is not very good at
estimating area or volume, so it turns out the estimate of
5000 tomatoes fitting in the truck is way off. That's why
people have a hard time at those contests where you are
supposed to estimate the number of jellybeans in a big jar.
Another example is that most people think their families use
about 10 gallons of water per day, but in reality the
average is about 300 gallons per day. When estimating area
or volume, you are much better off estimating linear
dimensions, and computing volume from the linear dimensions.
Here's a better solution to the problem about the tomato truck:

%\piccaption{ \qquad Albert Einstein, and his mustache, problem \ref{hw:hairmass}.}
\piccaption{ \qquad 
Can you guess how many jelly beans are in the jar? If you try to guess directly, you will almost certainly underestimate. The right way to do it is to estimate the linear dimensions, then get the volume indirectly.
See problem \ref{hw:jelly-beans}, p.~\pageref{hw:jelly-beans}.}
\parpic[r]{\includegraphics[width=0.3\textwidth]{./scaling/figs/jelly-beans_1.jpg} \label{jelly-beans}}


%\begin{figure}[h] 
%\begin{center}
%\includegraphics[width=0.4\textwidth]{./scaling/figs/jelly-beans}
%\caption{ \qquad 
%Can you guess how many jelly beans are in the jar? If you try to guess directly, you will almost certainly unde%restimate. The right way to do it is to estimate the linear dimensions, then get the volume indirectly.
%See problem \ref{hw:jelly-beans}, p.~\pageref{hw:jelly-beans}.}

%\end{center}
%\end{figure}

%\begin{figure}[h] 

\vspace{1.5mm}





\begin{exmp}
Cost of transporting tomatoes (correct solution): 
As in the previous solution, say the cost
of the trip is \$2000. The dimensions of the bin are
probably 4 m $\times$ 2 m $\times$ 1 m, for a volume of $8\ \munit^3$.
Since the whole thing is just an order-of-magnitude
estimate, let's round that off to the nearest power of ten,
$10\ \munit^3$. The shape of a tomato is complicated, and I don't
know any formula for the volume of a tomato shape, but since
this is just an estimate, let's pretend that a tomato is a
cube, 0.05 m $\times$ 0.05 m $\times$ 0.05 m, for a volume of $1.25\times10^{-4}\ \munit^3$.
Since this is just a rough estimate, let's round that to
$10^{-4} \munit^3$. We can find the total number of tomatoes by
dividing the volume of the bin by the volume of one tomato:
$10\ \munit^3/10^{-4}\ \munit^3=10^5$ tomatoes. The transportation cost
per tomato is $\$2000/10^5$ tomatoes=\$0.02/tomato. That
means that transportation really doesn't contribute very
much to the cost of a tomato.
\end{exmp}
\divider
\vspace{2mm}

Approximating the shape of a tomato as a cube is an example
of another general strategy for making order-of-magnitude
estimates. A similar situation would occur if you were
trying to estimate how many $\munit^2$ of leather could be
produced from a herd of ten thousand cattle. There is no
point in trying to take into account the shape of the cows'
bodies. A reasonable plan of attack might be to consider a
spherical cow. Probably a cow has roughly the same surface
area as a sphere with a radius of about 1 m, which would
be $4\pi (1\ \munit)^2$. Using the well-known facts that pi
equals three, and four times three equals about ten, we can
guess that a cow has a surface area of about $10\ \munit^2$, so
the herd as a whole might yield $10^5\ \munit^2$ of leather.

\piccaption{ \qquad Consider a spherical cow.}
\parpic[r]{\includegraphics[width=0.2\textwidth]{./scaling/figs/spherical-cow} \label{spherical-cow}}

\begin{exmp}\label{eg:amphicoelias}
Estimating Mass Indirectly: Usually the best way to estimate mass is to estimate linear dimensions, then use those to infer volume, and then get the mass based on the volume.
For example, \emph{Amphicoelias}, shown in the figure, may have been the largest
land animal ever to live. Fossils tell us the linear dimensions of an animal,
but we can only indirectly guess its mass. Given the length scale in the figure,
let's estimate the mass of an \emph{Amphicoelias}.

Its torso looks like it can be approximated by a rectangular box with dimensions
$10\ \munit\times5\ \munit\times3\ \munit$, giving about $2\times10^2\ \munit^3$. Living things
are mostly made of water, so we assume the animal to have the density of water,
$1\ \zu{g}/\zu{cm}^3$, which converts to $10^3\ \kgunit/\munit^3$. This
gives a mass of about $2\times10^5\ \kgunit$, or 200 metric tons.
\end{exmp}
\divider
\vspace{2mm}

\piccaption{ \qquad }
\parpic[r]{\includegraphics[width=0.4\textwidth]{./scaling/figs/amphicoelias}}






The following list summarizes the strategies for getting a
good order-of-magnitude estimate.

\begin{enumerate}
\item Don't even attempt more than one significant figure of precision.

\item Don't guess area, volume, or mass directly. Guess linear
dimensions and get area, volume, or mass from them.

\item When dealing with areas or volumes of objects with
complex shapes, idealize them as if they were some simpler
shape, a cube or a sphere, for example.

\item Check your final answer to see if it is reasonable. If
you estimate that a herd of ten thousand cattle would yield
$0.01\ \munit^2$ of leather, then you have probably made a mistake
with conversion factors somewhere.
\end{enumerate}



%%%%%%%%%%%%%%%%%%%%%%%%%%%%%%%%%%%%%%%%%%%%%%%%%%%%%%%%%


\section{Summary}

Nature behaves differently on large and small scales.
Galileo showed that this results fundamentally from the way
area and volume scale. Area scales as the second power of
length, $A\propto L^2$, while volume scales as length to the
third power, $V\propto L^3$.

An order of magnitude estimate is one in which we do not
attempt or expect an exact answer. The main reason why the
uninitiated have trouble with order-of-magnitude estimates
is that the human brain does not intuitively make accurate
estimates of area and volume. Estimates of area and volume
should be approached by first estimating linear dimensions,
which one's brain has a feel for.


\divider
\vspace{2mm}

\startexercises

\begin{enumerate}

\item  How many cubic inches are there in a cubic foot?
The answer is not 12.


\item 
Assume a dog's brain is twice as great in diameter as a
cat's, but each animal's brain cells are the same size and
their brains are the same shape. In addition to being a far
better companion and much nicer to come home to, how many
times more brain cells does a dog have than a cat?
The answer is not 2.

\item
 The population density of Los Angeles is about 
$4000\ \text{people}/\zu{km}^2$. That of San Francisco is about 
$6000\ \text{people}/\zu{km}^2$. How many times farther away is the average
person's nearest neighbor in LA than in San Francisco?
The answer is not 1.5. 


\item 
 A hunting dog's nose has about 10 square inches of active
surface. How is this possible, since the dog's nose is only
about 1 in $\times$ 1 in $\times$ 1 in = 1 $\zu{in}^3?$ After all, 10 is
greater than 1, so how can it fit?

\item 
 Estimate the number of blades of grass on a football field.

\item
 In a computer memory chip, each bit of information (a 0
or a 1) is stored in a single tiny circuit etched onto the
surface of a silicon chip. The circuits cover the surface of
the chip like lots in a housing development. A typical chip
stores 64 Mb (megabytes) of data, where a byte is 8 bits.
Estimate (a) the area of each circuit, and (b) its linear size.


\item  Suppose someone built a gigantic apartment building,
measuring 10 km $\times$ 10 km at the base. Estimate how tall the
building would have to be to have space in it for the entire
world's population to live.

\item 
 A hamburger chain advertises that it has sold 10 billion
Bongo Burgers. Estimate the total mass of feed required to
raise the cows used to make the burgers.

\item 
 Estimate the volume of a human body, in $\zu{cm}^3$.

\item 
  How many $\zu{cm}^2$ is 1 $\zu{mm}^2?$

\item 
  Compare the light-gathering powers of a 3-cm-diameter
telescope and a 30-cm telescope.

\item   One step on the Richter scale corresponds to a
factor of 100 in terms of the energy absorbed by something
on the surface of the Earth, e.g., a house. For instance, a
9.3-magnitude quake would release 100 times more energy than
an 8.3. The energy spreads out from the epicenter as a wave,
and for the sake of this problem we'll assume we're dealing
with seismic waves that spread out in three dimensions, so
that we can visualize them as hemispheres spreading out
under the surface of the earth. If a certain 7.6-magnitude
earthquake and a certain 5.6-magnitude earthquake produce
the same amount of vibration where I live, compare the
distances from my house to the two epicenters.

\item 
 In Europe, a piece of paper of the standard size,
called A4, is a little narrower and taller than its American
counterpart. The ratio of the height to the width is the
square root of 2, and this has some useful properties. For
instance, if you cut an A4 sheet from left to right, you get
two smaller sheets that have the same proportions. You can
even buy sheets of this smaller size, and they're called A5.
There is a whole series of sizes related in this way, all
with the same proportions. (a) Compare an A5 sheet to an A4
in terms of area and linear size. (b) The series of paper
sizes starts from an A0 sheet, which has an area of one
square meter. Suppose we had a series of boxes defined in a
similar way: the B0 box has a volume of one cubic meter, two
B1 boxes fit exactly inside an B0 box, and so on. What would
be the dimensions of a B0 box? 






\vspace{1.5mm}
\item 
 Estimate the mass of one of the hairs in Albert
Einstein's mustache, in units of kg.


\piccaption{ \qquad Albert Einstein, and his mustache, problem \ref{hw:hairmass}.}
\parpic[r]{\includegraphics[width=0.3\textwidth]{./scaling/figs/einstein} \label{hairmass}}


\item  According to folklore, every time you take a breath, you
are inhaling some of the atoms exhaled in Caesar's
last words. Is this true? If so, how many?

\item  The Earth's surface is about 70\% water. Mars's diameter
is about half the Earth's, but it has no surface water.
Compare the land areas of the two planets.

\item 
  The traditional Martini glass is shaped like a cone
with the point at the bottom. Suppose you make a Martini by
pouring vermouth into the glass to a depth of 3 cm, and then
adding gin to bring the depth to 6 cm. What are the
proportions of gin and vermouth?

\item 
 The central portion of a CD is taken up by the hole
and some surrounding clear plastic, and this area is
unavailable for storing data. The radius of the central
circle is about 35\% of the outer radius of the data-storing area.
What percentage of the CD's area is therefore lost? 





\newcommand{\tikzAngleOfLine}{\tikz@AngleOfLine}
  \def\tikz@AngleOfLine(#1)(#2)#3{%
  \pgfmathanglebetweenpoints{%
    \pgfpointanchor{#1}{center}}{%
    \pgfpointanchor{#2}{center}}
  \pgfmathsetmacro{#3}{\pgfmathresult}%
 }



\begin{tikzpicture}
    \coordinate (A) at (1,1);
    \coordinate (B) at ($(A)+(25:3)$);
    \coordinate (C) at ($(A)+(100:5)$);

\draw (A) node[left]{} -- (B) node[right]{}node[midway,below]{$C$} -- (C)node[above]{}node[midway,above]{$A$} -- (A)node[midway,left]{$B$};

    \tikzAngleOfLine(A)(B){\AngleStartAB}
    \tikzAngleOfLine(A)(C){\AngleEndAC}

    \tikzAngleOfLine(B)(A){\AngleStartBA}
    \tikzAngleOfLine(B)(C){\AngleEndBC}

    \tikzAngleOfLine(C)(A){\AngleStartCA}
    \tikzAngleOfLine(C)(B){\AngleEndCB}


    \draw[red,-] (A)+(\AngleStartAB:1cm) arc (\AngleStartAB:\AngleEndAC:1 cm);
    \node[circle] at ($(A)+({(\AngleStartAB+\AngleEndAC)/2}:0.5 cm)$) {$\alpha$};

    \draw[red,-] (B)+(\AngleStartBA:1cm) arc (\AngleStartBA:\AngleEndBC:1 cm);
    \node[circle] at ($(B)+({(\AngleStartBA+\AngleEndBC)/2}:0.5 cm)$) {$\beta$};

    \draw[red,-] (C)+(\AngleStartCA:1cm) arc (\AngleStartCA:\AngleEndCB:1 cm);
    \node[circle] at ($(C)+({(\AngleStartCA+\AngleEndCB)/2}:0.5 cm)$) {$\gamma$};

\end{tikzpicture}
\label{triangle-formula}




\item
The one-liter cube in the photo has been marked off into smaller
cubes, with linear dimensions one tenth those of the big one.
What is the volume of each of the small cubes?

\piccaption{ \qquad Problem \ref{hw:liter-cube}.}
\parpic[r]{\includegraphics[width=0.3\textwidth]{./scaling/figs/hw-liter-cube} \label{liter-cube}}






\vspace{1.5mm}

\item  
(a) Based on the definitions of the sine, cosine, and tangent, what units
must they have?
(b) A cute formula from trigonometry lets you find any angle of
a triangle if you know the lengths of its sides. Using the
notation shown in the figure, and letting $S=(A+B+C)/2$ be half the
perimeter, we have 
\begin{equation*}
 \tan \alpha/2=\sqrt{\frac{(S-B)(S-C)}{S(S-A)}} \qquad .
\end{equation*}
Show that the units of this equation make sense. In other words, check that the
units of the right-hand side are the same as your answer to part a of the question.

\item
Estimate the number of man-hours required for building the Great Wall of China. 

\item 
(a) Using the microscope photo in the figure, estimate the mass of a one cell of the
\emph{E. coli} bacterium, which is one of the most common ones in the human
intestine. Note the scale at the lower right corner, which is 1 $\mu\munit$. Each of
the tubular objects in the column is one cell.
(b) The feces in the human intestine are mostly bacteria (some dead, some alive),
of which \emph{E. coli} is a large and typical component. Estimate the number of bacteria
in your intestines, and compare with the number of human cells in your body, which is
believed to be roughly on the order of $10^{13}$. (c) Interpreting your result from part b,
what does this tell you about the size of a typical human cell compared to the size of
a typical bacterial cell?


\piccaption{ \qquad Problem \ref{hw:e-coli}.}
\parpic[r]{\includegraphics[width=0.4\textwidth]
{./scaling/figs/hw-e-coli} \label{e-coli}}


\item 
A taxon (plural taxa) is a group of living things. For example, \emph{Homo sapiens} and
\emph{Homo neanderthalensis} are both taxa --- specifically, they are two different species within
the genus \emph{Homo}. Surveys by botanists show that the number of plant taxa native to
a given contiguous land area $A$ is usually approximately proportional to $A^{1/3}$. (a) There are 70
different species of lupine native to Southern California, which has an area of about
$200,000\ \zu{km}^2$. The San Gabriel Mountains cover about $1,600\ \zu{km}^2$. Suppose
that you wanted to learn to identify all the species of lupine in the San Gabriels.
Approximately how many species would you have to familiarize yourself with?
(b) What is the interpretation of the fact that the exponent, $1/3$, is less than one?


\item 
X-ray images aren't only used with human subjects but also, for example,
on insects and flowers. In 2003, a team of researchers at Argonne National Laboratory
used x-ray imagery to find for the first time that insects, although they do not have
lungs, do not necessarily breathe completely passively, as had been believed previously; many insects rapidly compress and
expand their trachea, head, and thorax in order to force air in and out of their bodies.
One difference between x-raying a human and an insect is that if a medical x-ray machine was
used on an insect, virtually 100\% of the x-rays would pass through its body, and there would
be no contrast in the image produced. Less penetrating x-rays of lower energies have to
be used. For comparison, a typical human body mass is about 70 kg, whereas
a typical ant is about 10 mg. Estimate the ratio of the thicknesses of tissue that must be
penetrated by x-rays in one case compared to the other.

\item 
Radio was first commercialized around 1920, and ever since then, radio signals from our
planet have been spreading out across our galaxy. It is possible that alien civilizations
could detect these signals and learn that there is life on earth. In the 90 years that
the signals have been spreading at the speed of light, they have created a sphere with a radius of 90 light-years.
To show an idea of the size of this sphere, I've indicated it in the figure as a tiny white circle on an image
of a spiral galaxy seen edge on. (We don't have similar photos of our own Milky Way galaxy, because we can't
see it from the outside.) So far we haven't received answering signals from aliens within this sphere,
but as time goes on, the sphere will expand as suggested by the dashed outline, reaching more and more stars that might harbor extraterrestrial
life. Approximately what year will it be when the sphere has expanded to fill a volume 100 times greater than
the volume it fills today in 2010?


\piccaption{ \qquad Problem \ref{hw:seti-volume}.}
\parpic[r]{\includegraphics[width=0.4\textwidth]{./scaling/figs/seti-volume} \label{seti-volume}}



\vspace{1.5mm}

\item
Estimate the number of jellybeans in figure \ref{jelly-beans} on p.~\pageref{fig:jelly-beans}.


\item 
At the grocery store you will see oranges packed neatly in stacks.
Suppose we want to pack spheres as densely as possible, so that the greatest
possible fraction of the space is filled by the spheres themselves, not
by empty space. Let's call this fraction $f$.
Mathematicians have proved that the best possible result is $f\approx 0.7405$,
which requires a systematic pattern of stacking. If you buy ball bearings or golf
balls, however, the seller is probably not going to go to the trouble of stacking
them neatly. Instead they will probably pour the balls into a box and vibrate the
box vigorously for a while to make them settle. This results in a random packing.
The closest random packing has $f\approx 0.64$.
Suppose that golf balls, with a standard diameter of 4.27 cm, are sold in bulk
with the closest random packing. What is the diameter of the largest ball that
could be sold in boxes of the same size, packed systematically, so that there
would be the same number of balls per box?




\piccaption{ \qquad Problem \ref{hw:golf-ball-packing}.}
\parpic[r]{\includegraphics[width=0.24\textwidth]{./scaling/figs/oranges.jpg}\label{golf-ball-packing}}



\vspace{1.5mm}

\item 
Plutonium-239 is one of a small number of important long-lived forms of high-level
radioactive nuclear waste. The world's waste stockpiles have about $10^3$ metric tons
of plutonium. Drinking water is considered safe by U.S. government standards if it
contains less than $2\times10^{-13}\ \zu{g}/\zu{cm}^3$ of plutonium.
The amount of radioactivity to which you were exposed by drinking such water
on a daily basis would be very small compared to the natural background radiation that you are
exposed to every year. Suppose that the world's inventory of plutonium-239 were ground
up into an extremely fine dust and then dispersed over the world's oceans, thereby
becoming mixed uniformly into the world's water supplies over time.
Estimate the resulting concentration of plutonium, and compare with the government standard.

\end{enumerate}


  % need work
%%\documentclass[12pt]{book}
%\usepackage[width=5.5in,height=8.5in,
%  hmarginratio=3:2,vmarginratio=1:1]{geometry}
%
%% for some of these packages, you might have to install
%% texlive-latex-extra (in Ubuntu)
%
%\usepackage[T1]{fontenc}
%\usepackage{textcomp}
%\usepackage{mathpazo}
%\usepackage{pgfplots}
%%\pgfplotsset{compat=newest}
%%\usepackage{pslatex}
%
%\usepackage{url}
%\usepackage{fancyhdr}
%\usepackage{graphicx}
%\usepackage{subfig}
%\usepackage{amsmath}
%\usepackage{amsthm}
%%\usepackage{amssymb}
%\usepackage{makeidx}
%\usepackage{setspace}
%\usepackage{hevea}                           
%\usepackage{upquote}
%
%
%% these styles get translated in CSS for the HTML version
%\newstyle{a:link}{color:black;}
%\newstyle{p+p}{margin-top:1em;margin-bottom:1em}
%\newstyle{img}{border:0px}
%
%% change the arrows in the HTML version
%
%\makeindex
%
%
%\begin{document}


%
%
%\newcount\anchorcnt
%\newcommand*{\Anchor}[1]{%
%  \@bsphack%
%    \Hy@GlobalStepCount\anchorcnt%
%    \edef\@currentHref{anchor.\the\anchorcnt}% 
%    \Hy@raisedlink{\hyper@anchorstart{\@currentHref}\hyper@anchorend}% 
%    \M@gettitle{}\label{#1}% 
%    \@esphack%
%}
%
%
%
%%%% EXERCISE
%
%\newtheoremstyle{exercise}% name of the style to be used
%  {\topsep}% measure of space to leave above the theorem. E.g.: 3pt
%  {\topsep}% measure of space to leave below the theorem. E.g.: 3pt
%  {}% name of font to use in the body of the theorem
%  {0pt}% measure of space to indent
%  {\bfseries}% name of head font
%  {}% punctuation between head and body
%  { }% space after theorem head; " " = normal interword space
%  {}% Manually specify head
%
%\theoremstyle{exercise}
%\newtheorem{exercise}{Exercise}[chapter]
%
%%\newcounter{exercise}[chapter]
%%\newcommand{\nextexercise}{\refstepcounter{exercise}}
%
%%\newenvironment{exercise}{\nextexercise  \textbf{Exercise \thechapter.\theexercise} \begin{itshape} }{\end{itshape}}
%
%%\blankpage
%
%
%

\chapter{Descriptive statistics}
\label{descriptive}
\section{Significance}

Experiments are at the heart of the field of physics. The theories, laws and models relies on observations and experimentations. The gathering of data and analysis is an intrinsic part of physics. However, how do we go about turning data into theories, laws and models? Statistics is our friend  and in this chapter we will introduce a few notions about the analysis and interpretation of data. 

Say you want to determine the time of fall of steel balls in a cylinder filled with water over a distance of one meter.  You replicate the measurement several times  and collect your data into a table. One of your friend perform a similar experiment with a narrower cylinder. You find that on average the time of fall is slower by 2 ms in the narrow cylinder. 

\index{apparent effect}

A difference like that is called an {\bf apparent effect}; that is,
there might be something going on, but we are not yet sure.  There are
several questions we still want to ask:

\begin{itemize}

\item If the two experiments have different means, what about other {\bf
  summary statistics}, like median and variance?  Can we be more
  precise about how they differ?
\index{summary statistic}

\item Is it possible that the difference we saw could occur by chance,
  even if the experiments we compared were actually the same?  If so,
  we would conclude that the effect was not {\bf statistically
    significant}.
\index{statistically significant}
\index{significance}


\item Is it possible that the apparent effect is due to selection bias or
  some other error in the experimental setup?  If so, then we might
  conclude that the effect is an {\bf artifact}; that is, something we
  created (by accident) rather than found. 
\index{artifact}

\end{itemize}

Answering these questions will take most of the rest of this chapter.

\section{Errors}
 Before we go into details of error analysis it is important to understand the meaning of error in science. Error in a scientific measurement usually does not mean a mistake or blunder. Instead, the terms ``error'' and ``uncertainty'' both refer to unavoidable imprecision in measurements.

Of course, not all measurements have errors. If asked how many people there are in a room, one can usually give an exact number as an answer. However, if we want to know how many atoms there are in a room, giving an exact answer is nearly impossible, as the animation below illustrates.

In the physical science, you will deal mostly with measurements of the second kind. Since you will not be able to measure things with arbitrarily high precision, you should know how to quantify the imprecision of your results.
 Under different circumstances, the error that we are off by may vary greatly, so it is very useful to have a measurement of the quality of the measurement, which we call the {\bf error analysis}. First we choose a physical {\bf quantity} $X$ to determine and then we perform an {\bf experiment} which is a set of measurements of this quantity. Each {\bf measurement} gives a {\bf value} $x_i$. The {\bf outcome} $O_X$ of the experiment is the set of all results. Finally, we present the {\bf result} $\bar{x}$, which is the mean of the results in the outcome. The {\bf true value} $X_o$ is the value that would be measured if our experiment had no error. A good experiment will produce a result very close to the true value, $\bar{x} \cong X_o$.

Sources of error are commonly classified as either {\it systematic error} or {\it random error}. A {\bf systematic error} is any error that skews the measured values by the same amount in the same direction every time.\footnote{If an error source skews the measured values by a roughly constant amount in one direction, but has a slightly different effect each time, then it is probably an error source that has both a systematic and a random component.} For example, if you were making measurements with a stopwatch that always started at 0.01s instead of 0.00s, then you would have a systematic error. A {\bf random error} is any error that is not a systematic error. For example, if you were measuring an object with a ruler and you were off by a fraction of a millimeter, then you would have a random error. 

Random error usually skews measured values higher and lower with equal probability. If this is the case, then it possible to eliminate it by averaging, which we will denote averaging by enclosing the expression in angle brackets or covering it with a bar. First, we represent the value of a measurement by breaking it down into three components.

\[\textrm{Measurement = True Value + Systematic Error + Random Error}\]

%When we take the average of the measurements, we mean the exact average, which would be the average that you would obtain if you made the measurement an infinite number of times and averaged your results. 
Now averaging is distributive so we have

\[\langle \textrm{Measurement} \rangle  = \langle\textrm{True Value}\rangle + \langle\textrm{Systematic Error}\rangle + \langle\textrm{Random Error}\rangle \]

The first two terms on the right hand side are just constants, so we can drop the averaging symbols on them. The last term has an average of approximately zero in most cases yielding

\[\langle \textrm{Measurement} \rangle \textrm{ = True Value + Systematic Error}\]

We can't get rid of the systematic error with statistical methods, but by using quality equipment and improving our experimental technique we can usually reduce systematic error to a negligible level. Therefore, if you average your measured values and take care to minimize systematic error then you will obtain a result that is approximately equal to the true value.

It is not always the case that random error cancels with averaging. For example, if you are measuring the speed of a rolling ball and there is an intermittent breeze that is always in the same direction, then there will be a random error in the speed that biases it in one direction. Averaging will not be able to eliminate this type of random error. 

The most common measurements in the lab are done with devices that have a marked scale. There in an intrinsic limitation to the precision of our measurement. Say me measure the length of a pendulum. We find that is certainly closer to 128.9 cm than to 128.8 cm or 129.0 cm. Thus we can state with absolute confidence that the length is 128.8cm $\pm$ 0.1 cm. We call the first term, 128.9 cm, the "central value" and the second term, 0.1 cm, the "error" or "uncertainty" \footnote{We will revisit this subject later on.  Suffice to say this is a worst-case scenario. Most likely the error is much smaller. We will see later on that the standard uncertainty is given $\frac{(129.0-128.8)}{2\sqrt{6}}=0.04$~cm.}.

Many modern laboratory instruments use digital displays. For example, you can measure weights using an electronic scale. It has the advantage that you don't have to judge what is the central value of your measurement. It will give you the central value. All you have to do is write it down. Now, what is the error of this seemingly "precise" measurement? All instruments, no matter how sophisticated, have a limit to their precision. Usually, the manufacturer of the instrument will specify the error. (Most of the time it is written on the back of the instrument.) However, if you can't find this information, use the following rule of thumb: The error of an electronic device is usually half of the last precision digit. The following example should make it clear. Suppose you measured your weight with the electronic device pictured above. The result is 138.2 lbs. What is the error? If you can't find the manufacturer information about the error, you notice that the resolution of your measurement is 0.1 lb. In other words, the device claims that it can give you a more precise value than 138 lbs, but it doesn't tell you whether your weight is 138.15 lbs or 138.25 lbs. Thus we can assume that the error of our measurement is 0.05 lb. And our answer is 138.2 lbs $\pm$ 0.05 lb \footnote{Again this is an upper bound to the error,  the standard uncertainty is $\frac{(138.25-138.15)}{2\sqrt{3}}=0.03$~lb.}.

Not all measurements are done with instruments whose error can be reliably estimated. A classic example is the measuring of time intervals using a stopwatch. Of course, there will be a read-off error as discussed in the previous sections. However, that error will be negligible compared to the dominant error, the one coming from the fact that we, human beings, serve as the main measuring device in this case. Our individual reaction time in starting and stopping the watch will be by far the major source of imprecision. Since humans don't have built-in digital displays or markings, how do we estimate this dominant error?

The solution to this problem is to repeat the measurement many times. Then the average of our results is likely to be closer to the true value than a single measurement would be. 

\section{Means and averages}
\label{mean}

If you have a sample of $N$~values, $x_i$, the mean, $\bar{x}$, is
the sum of the values divided by the number of values; in other words
%
\[ \bar{x} = \frac{1}{N} \sum_{i=1}^N x_i \]
%

Another standard notation for the average is $\langle x\rangle$.

Now, suppose you measure the oscillation period of a pendulum with a stopwatch five times. You obtain the following data (in seconds):
\begin{verbatim}
3.55, 3.65, 3.73, 3.59, 3.67
\end{verbatim}

We find the mean period of oscillations is :

\[ \bar{x} = \frac{3.55+ 3.65+ 3.73+3.59+3.67}{5} = 3.65 \text{s}\]

The words ``mean'' and ``average'' are sometimes used interchangeably,
but we will maintain this distinction:

\begin{itemize}

\item The ``mean'' of a sample is the summary statistic computed with
  the previous formula.

\item An ``average'' is one of many summary statistics you might
  choose to describe the typical value or the
  {\bf central tendency} of a sample.
\index{central tendency}

\end{itemize}

Sometimes the mean is a good description of a set of values.  For
example, apples are all pretty much the same size (at least the ones
sold in supermarkets).  So if I buy 6 apples and the total weight is 3
pounds, it would be a reasonable summary to say they are about a half
pound each.
\index{weight!pumpkin}

But pumpkins are more diverse.  Suppose I grow several varieties in my
garden, and one day I harvest three decorative pumpkins that are 1
pound each, two pie pumpkins that are 3 pounds each, and one Atlantic
Giant that weighs 591 pounds.  The mean of
this sample is 100 pounds, but if I told you ``The average pumpkin
in my garden is 100 pounds,'' that would be wrong, or at least
misleading.
\index{pumpkin}

In this example, there is no meaningful average because
there is no typical pumpkin.



\section{Distributions}
\label{distributions}
\index{distribution}

Summary statistics are concise, but dangerous, because they obscure
the data.  An alternative is to look at the {\bf distribution} of the
data, which describes how often each value appears.

The most common representation of a distribution is a {\bf histogram},
which is a graph that shows the frequency or probability
of each value.\index{histogram}. A {\bf histogram} is a chart made by stacking up squares in columns based on the values of measurements, with one square added per measurement. This makes the more frequent results form taller columns.

\begin{figure}[h]
\begin{center}
\begin{tikzpicture}
\begin{axis}[
ybar interval,
xtick={1,2,3,4,5,6,7,8},% reset from ybar interval
%xticklabel=
%{$[\pgfmathprintnumber\tick,%
%\pgfmathprintnumber\nexttick)$}
%]
]
% a data file containing 8000 normally distributed
% random numbers of mean 0 and variance 1

\addplot coordinates {
    (1,1) (2,5) (3,16)
     (4,26) (5,12) (6,4) (7,5) (8,2)};
\end{axis}
\end{tikzpicture}
\caption{Sample Histogram}
\end{center}
\end{figure}

In this figure, each square represents one measurement of a quantity $X$. The value of the measurement, which we call $x$, is the number of the column that it lies in. We can write the value of the $i^{th}$ measurement as $x_i$. Most quantities that we measure are not always integers, but we can still use histograms by making each column a {\bf bin} that collects any value falling within a specific range. We will refer to the shape that the histogram forms as a {\bf distribution} of measurements. 

There are two distinct measures of the error of a result, {\it accuracy} and {\it precision}. The {\bf accuracy} tells you how close your result is to the accepted result. Accuracy is almost always expressed in terms of the {\bf percent error}
\footnote{Percent error is also called percent deviation}, which is given by 
\begin{equation}\frac{\textrm{Experimental Result} - \textrm{Accepted Result}}{\textrm{Accepted Result}} \times 100 \end{equation}

The {\bf precision} tells you how consistent your measured values are, with higher precision indicating that all of your data points are clustered around one value. Effectively, precision is a measure of the width of the distribution for the experiment. In order to quantify the width of the distribution we need a formula that accounts for the contribution to the spread by each measured value. The first thing that comes to mind is the {\bf deviation} of each value, which is given by 
\begin{equation} d_i \equiv x_i - \bar{x} \end{equation}

but this won't help to much because if we average the deviation over all the values, we will usually get an answer close to zero. This is because random error usually is equally likely to skew the values in either direction from the average and this will cause the deviations to cancel positive with negative. But there is a quick fix for this---we can just take the absolute value of the deviation and average that. This defines the {\bf mean deviation} \footnote{The mean deviation is also called the average deviation}
\begin{equation} \alpha_{O_X} = \frac{1}{N} \sum_i |x_i - \bar{x}| \end{equation}
where the subscript $O_X$ indicates that the expression is computed using the outcome of the experiment measuring the quantity $X$.

The mean deviation is not used very often. Instead of taking the absolute value, the square of the deviation is used. The square is again always positive, but it is more convenient than the absolute values because squares are easier to work with analytically and there are other advantages that we will see soon. Averaging the square of the deviation gives us the {\bf variance} 
\footnote{This is the biased or population standard deviation. There are arguments that it is better to divide by N-1 instead of N in the average, but we will not pursue that here}
\begin{equation}\sigma^2_{O_X} = \frac{1}{N} \sum_{i=1}^N \left(x_i - \bar{x}\right)^2  \end{equation}
%
%There is a useful identity involving the standard deviation that can be derived easily.
%\[\sigma_{O_X}^2 = \left\langle (x_i - \langle x_i \rangle)^2 \right\rangle 
%= \left\langle x_i^2 - 2x_i\langle x_i \rangle + \langle x_i \rangle^2 \right\rangle \]
%\[ = \langle x_i^2 \rangle - \langle 2x_i\langle x_i \rangle \rangle + \langle x_i \rangle^2 \]
%\[ = \langle x_i^2 \rangle - 2\langle x_i \rangle^2 + \langle x_i \rangle^2 
% = \langle x_i^2 \rangle - \langle x_i \rangle^2\]

In the same way that the mean is intended to describe the central
tendency, variance is intended to describe the {\bf spread}.
The variance of a set of values is\footnote{You may have seen or learned another definition of the variance in which $N-1$ is replaced by $N$. Technically the variance must include the $N$ factor whereas the sample variance includes the $N-1$ factor. The sample variance is an ``unbias'' estimator of the variance. It's more of a mathematical subtlety, which does not affect our reasoning here. In any case, the difference between the two is negligible for large $N$. See
\url{http://wikipedia.org/wiki/Bias_of_an_estimator} for more information.
 }
%
\[ s^2 = \frac{1}{N-1} \sum_{i=1}^N \left(x_i - \bar{x}\right)^2 \]
%
The term $x_i-\bar{x}$~is called the ``deviation from the mean,'' so
variance is the mean squared deviation, which is why it is denoted $s^2$.  


 
The variance is our previous example is
\begin{align*}
 s^2 =\frac{1}{5-1} [&(3.55-3.65)^2+ (3.65-3.65)^2+ (3.73-3.65)^2\\
 &+(3.59-3.65)^2+(3.67-3.65)^2 ] = 0.0051  \text{s}^2
\end{align*}




The square root of variance, $s$, is called the {\bf
  standard deviation}.
\index{deviation}
\index{standard deviation}


By itself, variance is hard to interpret.  One problem is that the
units are strange; in this case the measurements are in pounds, so the
variance is in pounds squared.  Standard deviation is more meaningful;
in this case the units are pounds.



Now we want to consider the shape that a histogram of the outcome would take if we had performed an infinite number of measurements. This limiting distribution would give us a sense for what the error itself really looks like instead of just seeing a few pieces of the puzzle. Without further information about the measurement process, we can't be sure what the histogram would look like. But the {\bf central limit theorem} proves that if the measurement process is affected by many independent random errors, then the limiting distribution of an infinite number of measurements will be approximately a {\it Gaussian distribution}. Whether or not this is the reason, it is true that very often random error in measurements produces a {\bf Gaussian distribution}, which is a bell-shaped curve with the functional form
\begin{equation} f(x) = \frac{1}{\sigma \sqrt{2 \pi}}  e^{-\frac{(x-\mu)^2}{2\sigma^2}}\end{equation}
where $\sigma$ is the width parameter, and $\mu$ is the center of the peak. These parameters are just numbers to plug into the formula, but their symbols were chosen because $\sigma$ is the standard deviation of the limiting distribution and $\mu$ is the average of the limiting distribution.

\begin{figure}[h]
\begin{center}
\definecolor{qqqqff}{rgb}{0,0,1}
\definecolor{qqqqqq}{rgb}{1,0,0}
%\begin{tikzpicture}[line cap=round,line join=round,>=triangle 45,x=1.0cm,y=10.0cm]
\begin{tikzpicture}[line cap=round,line join=round,>=,x=1.0cm,y=10.0cm]
\draw[->,color=black] (-3.99623,0) -- (3.81526,0);
\draw[->,color=black] (-3.5,-0.08169) -- (-3.5,0.49561);
%\clip(-3.49623,-0.38169) rectangle (3.81526,0.69561);
\draw[color=qqqqff,fill=qqqqff,fill opacity=0.25, smooth,samples=50,domain=-1.00:1.00] plot(\x,{1/sqrt(2*3.1415926535)*exp(-0.5*(\x)^2)}) -- (1.0,0) -- (-1.0,0) -- cycle;
\draw[line width=2pt,color=qqqqqq, smooth,samples=100,domain=-3.4962341176758:3.81525635549043] plot(\x,{1/sqrt(2*3.1415926535)*exp(-0.5*(\x)^2)});
\draw [color=black](1.00,-0.02) node[anchor=north] {$\mu+\sigma$};
\draw [color=black](-1.00,-0.02) node[anchor=north] {$\mu-\sigma$};
\draw [color=yellow](-0.54691,0.1499) node[anchor=north west] {\textbf{Area}};
\draw [color=yellow](-0.54691,0.0999) node[anchor=north west] {\textbf{0.68}};
\draw[color=blue] (-4.2325,0.45082) node {$y$};
\draw[color=blue] (3.23,-0.0582) node {$x$};
\draw[color=blue] (-0.03,-0.0582) node {$\mu$};
\fill [color=black] (1.00,0) circle (1.5pt);
\fill [color=black] (-1.00,0) circle (1.5pt);
\end{tikzpicture}

\caption{Gaussian Distribution: $ f(x) = \frac{1}{\sigma \sqrt{2 \pi}}  e^{-\frac{(x-\mu)^2}{2\sigma^2}}$. The probability that the value $x$ ranges from $\mu -\sigma$  to $\mu +\sigma$ is 68\%.  }
\end{center}
\end{figure}

The gaussian distribution is an example of a probability density function (PDF). Evaluating a PDF for a particular value of $x$~is usually not useful.
The result is not a probability; it is a probability {\em density}.
\index{density}
\index{mass}

In physics, density is mass per unit of
volume; in order to get a mass, you have to multiply by volume or,
if the density is not constant, you have to integrate over volume.
\index{inertia}
\index{moment of inertia}

% Similarly, probability density measures probability per unit of $x$.
% In order to get a probability mass, you have to integrate over $x$.
% For example, if $x$~is a random variable whose PDF is PDF$_X$, we
% can compute the probability that a value from $x$~falls between 
% $-$0.5 and 0.5:

% \[ P(-0.5 \le X < 0.5) = \int_{-0.5}^{0.5} PDF_{X}(x) dx \]


In this context, {\bf frequency} means the number of times a value
appears in a dataset---it has nothing to do with the pitch of a sound
or tuning of a radio signal.  A {\bf probability} is a frequency expressed
as a fraction of the sample size, $N$.
\index{frequency}
\index{probability}
\index{dictionary}


The result is a dictionary that maps from values to frequencies.
To get from frequencies to probabilities, we divide through by $N$,
which is called {\bf normalization}:\index{normalization}

\begin{equation}
\text{Relative frequency} = \text{the number of readings in the bin}
\end{equation}

The normalized histogram is called a {\bf PF}, which stands for
``probability function''; that is, it's a function that maps from
values to probabilities.
\index{PF}
\index{probability function}

To illustrate the idea, suppose we measure the time of fall of the steel balls  as discussed in the previous section, until we have 12 measurements (in seconds):

2.3, 2.4, 2.5, 2.6, 2.5, 2.3, 2.4, 2.4, 2.2, 2.4, 2.5, 2.5

Histograms and probability functions are most often plotted as bar charts.  Figure~\ref{pendulum} shows a histogram of the oscillation period of pendulum. 

\begin{figure}[h]
\begin{center}
\begin{tikzpicture}
  \centering
  \begin{axis}[
        ybar,
    height=6cm,
    width=10cm,
    enlarge y limits=false,
    axis lines*=left,
    ymin=0,
    ymax=5,
     legend style={at={(0.5,-0.2)},
        anchor=north,legend columns=-1},
        ylabel={Occurences},
        xlabel={Period (s)},
     xtick=data,
        nodes near coords,
    ]
    \addplot coordinates {
    (2.2,1) (2.3,2) (2.4,4)
      (2.5,3) (2.6,2)};
  \end{axis}
\end{tikzpicture}
\end{center}
\caption{ \quad Histogram of the period of pendulum.}
\label{pendulum}
\end{figure}

Histograms are useful because they make the following features immediately
apparent:

\begin{description}

\item[Mode:] The most common value in a distribution is called the
  {\bf mode}.  In Figure~\ref{pendulum} there is a clear mode at 2.4 s.  In this case, the mode is the summary statistic that does
  the best job of describing the typical value.
\index{mode}

\item[Shape:] Around the mode, the distribution is asymmetric; it
  drops off quickly to the left and more slowly to the right.  In statistics this is often referred to as a skew. 
\index{shape}

\item[Outliers:] Values far from the mode are called {\bf outliers}.
  Many of them are probably due to errors, either in the reporting
  or recording of data.
\index{outlier}

\end{description}

Although histograms make some features apparent, they are usually not
useful for comparing two distributions.  We can address this problem using probability functions.


\section{Probability Functions (PF's)}
\index{plotting}
\index{PF}


To plot a probability function as a bar graph:  Bar graphs are most useful if the number
of values in the probability function is small.

\begin{figure}[h]
\begin{tikzpicture}
  \centering
  \begin{axis}[
        ybar,
    height=6cm,
    width=13cm,
    enlarge y limits=false,
    axis lines*=left,
    ymin=0,
    ymax=0.3,
     legend style={at={(0.5,-0.4)},
        anchor=north,legend columns=-1},
        ylabel={Frequency},
        xlabel={Fall Time (s)},
     xtick=data,
        nodes near coords,
    every node near coord/.append style={
        anchor=mid west,
        rotate=70
    }
    ]
    \addplot coordinates { (3.3, 0.01) (3.4, 0.04) (3.5, 0.11) (3.6, 0.206) (3.7, 0.252)(3.8, 0.20) (3.9, 0.113)  (4.0, 0.0417)  };
        \addplot coordinates { (3.3, 0.015) (3.4, 0.05) (3.5, 0.13) (3.6, 0.236) (3.7, 0.242)(3.8, 0.18) (3.9, 0.103)  (4.0, 0.0317)   };
    \legend{wide cylinder, narrow cylinder}
  \end{axis}
\end{tikzpicture}
\caption{ \quad probability function of the fall time for narrow and wide cylinders}
\label{narrow_wide_cylinders}
\end{figure}

Figure~\ref{narrow_wide_cylinders} shows the PDF of our two experiments involving dropping steel balls in water cylinders as a bar
graph.  Using the probability function, we can see more clearly where the distributions
differ.  The steel balls in the narrow cylinders take slightly less time to fall than the wide cylinders.


\index{PF object}

probability function and histograms objects are similar in many ways.   The
biggest difference is that a histograms maps from values to integer
counters; a probability function maps from values to probabilities.

%
\index{mean}
\index{variance}
\index{PDF}
%
In Section~\ref{mean} we computed the mean of a sample by adding up
the elements and dividing by $N$.  If you are given a probability function, you can
still compute the mean, but the process is slightly different:
%
\[ \bar{x} \ = \sum_i p_i x_i \]
%
where the $x_i$~are the unique values in the probability function and $p_i$=PF($x_i$).
Similarly, you can compute variance like this:
%
\[ s^2 = \sum_i p_i (x_i - \bar{x})^2\]



The standard deviation may serve as a measure of uncertainty. 
In physical science, for example, the reported standard deviation of a group of repeated measurements gives the precision of those measurements. When deciding whether measurements agree with a theoretical prediction, the standard deviation of those measurements is of crucial importance: if the mean of the measurements is too far away from the prediction (with the distance measured in standard deviations), then the theory being tested probably needs to be revised. This makes sense since they fall outside the range of values that could reasonably be expected to occur, if the prediction were correct and the standard deviation appropriately quantified. 


When presenting results it is desirable to include an estimate of the error right with the result itself. This is accomplished by appending an interval that the true value should lie within. By convention, the interval used is just the standard deviation of the mean. It is given by

\[ \sigma_{mean} = \frac{1}{\sqrt{N}} s \]

 It is true that there is a reasonable chance that the true value will not fall within one standard deviation of the result, but under normal circumstances the only distance that we could be sure the true value lies within would be far too large to be of use. This is because there is usually a slight probability that any given measurement will give a value that is very far from the true value. Therefore it is convenient to just state results as \begin{equation} X_o = \bar{x} \pm \sigma_{mean X} \end{equation}

After this form is obtained, it can be compared with other experiments for {\it consistency}. Two results $X_o = \bar{x} \pm \sigma_{mean X}$ and $X_o' = \bar{x}' \pm \sigma_{mean X}'$ are {\bf consistent} if  $|\bar{x} - \bar{x}'| \leq \sigma_{mean X} + \sigma_{mean X}'$.\footnote{This definition can be relaxed a little because of the fact that there is a decent chance that the true value lies outside one standard deviation in both experiments.}
In other words, two results are consistent if and only if the half width at half maximum of their probability distributions overlap. 



\section{Worst-Case Error Propagation}
So far we have seen the standard deviation used on the outcome of an experiment to get a sense of the spread of the data. It is also possible to estimate the expected standard deviation of the outcome by accounting for each of the random error sources, which is called {\bf error propagation}. It will take some mathematical work to derive the error propagation formula, but first we need to know how to account for the random error sources. First of all we must assume that we are aware of all the random error sources. This is usually not a problem because we just treat each measured quantity as having one combined random error source, which may in fact include many independent error sources. If we are directly measuring the quantity of interest, then that is all there is to it---the expected standard deviation is the standard deviation of the distribution of measurements. The only problem is that we do not know the distribution of measurements before we make them, so how can we know it's standard deviation? This is where we can make an estimate. The {\bf least count} is the size of the smallest division on the measuring device. A competent experimenter using working equipment should rarely make errors larger than the least count. In fact, you should always be able to round to the correct division, so you will probably always be within half of the least count. For example, on a normal ruler the least count is 1mm and you should never be off by more than 0.5mm. But clearly there is a good chance that your measured value will differ from the true value by 0.1mm, so there is some nonzero spread to the distribution of measurements. A good rule to use is to assume the standard deviation of your distribution of measurements is half of the least count. This means that about $68\%$ of your measurements lie within half a least count from the true value. Note that this asssumes that the systematic error has been brought down to a negligible level.

So we have handled the case where we are directly measuring the quantity of interest, but that is not usually how experiments are done. Most experiments involve several measurements whose results are plugged into an equation to get the final result. For example, if you are measuring the area of a rectangle, you must measure the width and the height, and then plug these results into the formula $A = w*h$. The error in measuring the width and height is known based on the least count, but what amount of error do these error sources produce in the result for the area? This is where the error propagation formulas are needed.

There is only one error propagation formula that you need to know, which will be derived at the end, but there are other formulas which are either approximations or special cases of the final formula. The simplest type of error propagation formula is the worst case type. Formulas of this type just find the highest and lowest possible values of the result for values of the parameters which lie within their error bounds. The formulas for addition and multiplication of parameters are shown below.

\[ F(X,Y) = X + Y \Rightarrow \sigma_F = \sqrt{\sigma^2_X + \sigma^2_Y} \]
\[ F(X,Y) = X*Y \Rightarrow \frac{\sigma_F}{|\bar{f}|} =\sqrt{ \left(\frac{\sigma_X}{|\bar{x}|}\right)^2 + \left(\frac{\sigma_Y}{|\bar{y}|}\right)^2} \]

Error in the form $\sigma_X/\bar{x}$ is called {\bf relative error}. In contrast, error in the regular form $\sigma_X$ is called {\bf absolute error}.
We will not be using this simple form of error propagation because it usually gives an overestimate of the error because it is not so likely that the results of both parameters will be off from the true value in the same direction.



\section{Standard scores}

In this section we look at relationships between variables.  For
example, we have a sense that height is related to weight; people who
are taller tend to be heavier.  {\bf Correlation} is a description of
this kind of relationship.
\index{correlation}

A challenge in measuring correlation is that the variables we want
to compare might not be expressed in the same units.  For example, height
might be in centimeters and weight in kilograms.  And even if they are
in the same units, they come from different distributions.
\index{units}



The standard solution to this problem is to transform all values to {\bf standard scores}.  This leads to
the Pearson coefficient of correlation.
\index{standard scores}
\index{Pearson coefficient of correlation}
\index{coefficient!correlation}




If $X$~is a series of values, $x_i$, we can convert to standard
scores by subtracting the mean and dividing by the standard deviation:
\[ z_{i}~=\frac{x_{i}-\mu}{\sigma}. \]
\index{mean}
\index{standard deviation}

The numerator is a deviation: the distance from the mean.  Dividing by
$\sigma$~{\bf normalizes} the deviation, so the values of $Z$~are
dimensionless (no units) and their distribution has mean 0 and
variance 1.
The absolute value of $z_i$ represents the distance between the raw score and the population mean in units of the standard deviation.
$z$ is negative when the raw score is below the mean, positive when above.

\index{normalize}
\index{deviation}
\index{normal distribution}
\index{distribution!normal}
\index{Gaussian distribution}
\index{distribution!Gaussian}

\section{Covariance}
\index{covariance}
\index{deviation}

{\bf Covariance} is a measure of the tendency of two variables
to vary together.  If we have two series, $X$~and $Y$, their
deviations from the mean are

\[ \Delta x_{i} = x_i-\mu_X \]

\[ \Delta y_{i} = y_i-\mu_Y \]

where $\mu_{X}$ is the mean of $X$~and $\mu_{Y}$ is the mean of $Y$.
If $X$~and $Y$~vary together, their deviations tend to have the same
sign.

If we multiply them together, the product is positive when the
deviations have the same sign and negative when they have the opposite
sign.  So adding up the products gives a measure of the tendency to
vary together.

Covariance is the mean of these products:
%
\[ Cov(X,Y) = \frac{1}{n} \sum_{i}^{n} \Delta x_i  \Delta y_i \]
%
where $n$~is the length of the two series (they have to be the same
length).

Covariance is useful in some computations, but
it is seldom reported as a summary statistic because it is hard to
interpret.  Among other problems, its units are the product of the
units of $X$~and $Y$.  So the covariance of weight and height might be
in units of kilogram-meters, which doesn't mean much.


It is easy to convinve yourself that the covariance of a list with itself is the variance
Var($X$)=Cov($X$, $X$).


\section{Correlation}
\index{correlation}
\index{standard score}

Since covariance is a dimensionful quantity, its value depends on the dimensions of $X$ and $Y$. 
One solution to this problem is to divide the deviations by the standard deviation $\sigma$,
which yields standard scores, and compute the product of standard scores:
%
\[ p_i = \frac{(x_i - \mu_X)}{\sigma_X} \frac{(y_i - \mu_Y)}{\sigma_Y} \]
%
The mean of these products is
%
\[ \rho = \frac{1}{n} \sum_i^n p_i \]
%
This value is called {\bf Pearson's correlation} after Karl Pearson,
an influential early statistician.  It is easy to compute and easy to
interpret.  Because standard scores are dimensionless, so is $\rho$.
\index{Pearson, Karl}
\index{Pearson coefficient of correlation}
\index{coefficient!correlation}

Also, the result is necessarily between $-1$ and +1.  To see why, we
can rewrite $\rho$~by factoring out $\sigma_{X}$ and $\sigma_Y$:
%
\[ \rho = \frac{Cov(X,Y)}{\sigma_X \sigma_Y} \]
%
Expressed in terms of deviations, we have
%
\[ \rho = \frac{\sum \Delta x_i \Delta y_x}{\sqrt{\sum (\Delta x_i)^2 \sum (\Delta y_i)^2}} \]
%
Then by the ever-useful Cauchy-Schwarz inequality\footnote{See
  \url{http://wikipedia.org/wiki/Cauchy-Schwarz_inequality}.}, we can show
that $\rho^{2} \le 1$, so $-1~\le~\rho~\le~1$.
\index{Cauchy-Schwarz inequality}

The magnitude of $\rho$~indicates the strength of the correlation.  If
$\rho~=~1$ the variables are perfectly correlated, which means that if
you know one, you can make a perfect prediction about the other.  The
same is also true if $\rho~=~-1$.  It means that the variables
are negatively correlated, but for purposes of prediction, a
negative correlation is just as good as a positive one.
\index{prediction}

Most correlation in the real world is not perfect, but it
is still useful.  For example, if you know someone's height, you might
be able to guess their weight.  You might not get it exactly right, but
your guess will be better than if you didn't know the height.
Pearson's correlation is a measure of how much better.

So if $\rho~=~0$, does that mean there is no
relationship between the variables?  Unfortunately, no.  Pearson's
correlation only measures {\em linear} relationships.  If there's a
nonlinear relationship, $\rho$~understates the strength of the
dependence.
\index{linear relationship}

\begin{figure}
% descriptive.py
\centerline{\includegraphics[height=2.5in]{statistics/Correlation_examples.png}}
\caption{Examples of datasets with a range of correlations.}
\label{corr_examples}
\end{figure}

Figure~\ref{corr_examples} is from
\url{http://wikipedia.org/wiki/Correlation_and_dependence}.  It shows
scatterplots and correlation coefficients for several
carefully-constructed datasets.
\index{coefficient!correlation}
\index{scatter plot}
\index{plot!scatter}

The top row shows linear relationships with a range of correlations;
you can use this row to get a sense of what different values of
$\rho$~look like.  The second row shows perfect correlations with a
range of slopes, which demonstrates that correlation is unrelated to
slope (we'll talk about estimating slope soon).  The third row shows
variables that are clearly related, but because the relationship is
non-linear, the correlation coefficient is 0.

The moral of this story is that you should always look at a scatterplot of
your data before blindly computing a correlation coefficient.
\index{correlation coefficient}
\index{coefficient!correlation}

\section{Least squares fit}

Correlation coefficients measure the strength and sign of a
relationship, but not the slope.  There are several ways to estimate
the slope; the most common is a {\bf linear least squares fit}.  A
``linear fit'' is a line intended to model the relationship between
variables.  A ``least squares'' fit is one that minimizes the mean
squared error (MSE) between the line and the data\footnote{See
  \url{http://wikipedia.org/wiki/Simple_linear_regression}.}.
\index{least squares fit}
\index{linear least squares}
\index{coefficient!correlation}
\index{linear regression}

Suppose we have a sequence of points, $Y$, that we want to express as a
function of another sequence $X$.  If there is a linear relationship
between $X$~and $Y$~with intercept $b$~and slope $m$, we
expect each $y_{i}$ to be roughly $m x_i + b$.
\index{residual}

But unless the correlation is perfect, this prediction is only
approximate.  The deviation, or {\bf residual}, is 

\[ \epsilon_{i}~=~(m x_{i}+b)~-~y_{i} \]

The residual might be due to random factors like measurement error,
or non-random factors that are unknown.  For example, if we are
trying to predict weight as a function of height, unknown factors
might include diet, exercise, and body type.
\index{slope}
\index{intercept}

If we get the parameters $b$~and $m$~wrong, the residuals
get bigger, so it makes intuitive sense that the parameters we want
are the ones that minimize the residuals.

As usual, we could minimize the absolute value of the
residuals, or their squares, or their cubes, etc.  The most common
choice is to minimize the sum of squared residuals
%
\[ \min_{m,b} \sum \epsilon_i^2 \]
%
Why?  There are three good reasons and one bad one:

\begin{itemize}

\item Squaring has the obvious feature of treating positive and
negative residuals the same, which is usually what we want.

\item Squaring gives more weight to large residuals, but not
so much weight that the largest residual always dominates.

\item If the residuals are independent of $x$, random, and normally
  distributed with $\mu~=~0$ and constant (but unknown) $\sigma$, then
  the least squares fit is also the maximum likelihood estimator of
  $b$~and $m$.\footnote{See Press et al., {\em Numerical Recipes in C},
    Chapter 15 at \url{http://www.nrbook.com/a/bookcpdf/c15-1.pdf}.}
\index{MLE}
\index{maximum likelihood estimator}

\item The values of $\hat{m}$ and $\hat{b}$ that minimize
  the squared residuals can be computed efficiently.

\end{itemize}

That last reason made sense when computational efficiency was more
important than choosing the method most appropriate to the problem
at hand.  That's no longer the case, so it is worth considering
whether squared residuals are the right thing to minimize.
\index{computation}

For example, if you are using values of $X$~to predict values of $Y$,
guessing too high might be better (or worse) than guessing too low.
In that case you might want to compute some cost function,
cost $(\epsilon_{i})$, and minimize total cost.
\index{cost function}

However, computing a least squares fit is quick, easy and often good
enough, so here's how:

\begin{enumerate}

\item Compute the sample means, $\bar{x}$~and $\bar{y}$, the variance
of $X$, and the covariance of $X$~and $Y$.

\item The estimated slope is
%
\[ \hat{m} = \frac{Cov(X,Y)}{Var(X)} \]
%
\item And the intercept is
%
\[ \hat{b} = \bar{y} - \hat{m} \bar{x} \]
%
\end{enumerate}

To see how this is derived, you can read
\url{http://wikipedia.org/wiki/Numerical_methods_for_linear_least_squares}.




\section{Goodness of fit}
\index{goodness of fit}
\index{fit, goodness}

Having fit a linear model to the data, we might want to know how good
it is.  Well, that depends on what it's for.  One way to evaluate a
model is its predictive power.

In the context of prediction, the quantity we are trying to guess is
called a {\bf dependent variable} and the quantity we are using to
make the guess is called an {\bf explanatory} or {\bf independent
  variable}.
\index{dependent variable}
\index{explanatory variable}
\index{independent variable}
\index{coefficient!determination}

To measure the predictive power of a model, we can compute the {\bf
  coefficient of determination}, more commonly known as ``R-squared'':
%
\[ R^2 = 1 - \frac{Var(\epsilon)}{Var(Y)}\]
%
To understand what $R^2$ means, suppose (again) that you are trying
to guess someone's weight.  If you didn't know anything about them,
your best strategy would be to guess the average weight, $\bar{y}$; in
that case the error of your guesses would be $Var(Y)$:
%
\[\frac{1}{n} \sum (\bar{y} - y_i)^2 = Var(Y) \]
%
But if I told you their height $x_i$, you would guess $\hat{m} x_i + \hat{b} $; in that case your mean standard error (MSE) would be $Var(\epsilon)$.
%
\[ MSE = 
\frac{1}{n} \sum ((\hat{m} x_i + \hat{b}) - y_i)^2 =
Var(\epsilon) \]
%
So the term $Var(\epsilon)/Var(Y)$ is the ratio of mean squared error with
and without the explanatory variable, which is the fraction of
variability left unexplained by the model.  The complement, $R^2$,
is the fraction of variability explained by the model.
\index{variability}

If a model yields $R^2~=~0.64$, you could say that the model explains
64\% of the variability, or it might be more precise to say that it
reduces the MSE of your predictions by 64\%. In the context of a linear least squares model, it turns out that
there is a simple relationship between the coefficient of
determination and Pearson's correlation coefficient, $\rho$:
\[ R^2 = \rho^2  \]

$R^2$ is often refer to as the correlation coefficient.

In practice a calculator or computer is used to calculate the values $m, b$ and $R^2$.  For reference, in EXCEL, the regression coefficients can be found as

\[b={\rm INTERCEPT}\left(Y_{ARRAY},X_{ARRAY}\right)\] 
\[m={\rm SLOPE}(Y_{ARRAY},X_{ARRAY})\] 
\[R^2={\rm RSQ}(Y_{ARRAY},X_{ARRAY})\]
 
It turns out that Excel has a particularly convenient utility for carrying out such calculations: A function called LINEST (which stands for LINE STatistics). To access the uncertainties in the slope and intercept do the following:

Using the mouse, highlight a region of empty cells that is two columns wide by five rows tall. This is where you will be inserting the results from LINEST.
2. Without deselecting the highlighted region, go to the formula bar and type the following "=LINEST(ydata,xdata,true,true)" where "ydata" indicates the cell numbers corresponding to the range of y data (e.g. it might read B2:B10 if you have y data in cells B2, B3,...,B10). The words "true" ought to literally appear as true comma true. Next, instead of hitting "Enter" simultaneously type 
"shift-ctrl-enter". That is, depress and hold down the shift key, the
control key, and finally, the enter key. A $2\times 5$ matrix of numbers should appear in the highlighted cells. The upper left column is the
slope, while the upper right number is the intercept.
The second row holds the uncertainties in those numbers. In particular,
the uncertainty in the slope is the second number down from the top in the left-most column.



\section{Error propagation rules}

\begin{itemize}
\item
The {\em Absolute Error} of a quantity $Z$ is given by $\sigma(Z)$, is equal or larger than zero.
\item
The {\em Relative Error} of a quantity $Z$ is given by $\frac{\sigma(Z)}{|Z|}$, is equal or larger than zero.
\item
To determine the error in a quantity $Z$ that is the {\em sum} of other quantities,
you {\em add the absolute errors} of those quantities (Rules 2,3 below). 
To determine the error in a quantity $Z$ that is the {\em product} of other quantities,
you {\em add the relative errors} of those quantities (Rules 4,5 below).
\end{itemize}

%%\begin{table}[htp]
\begin{center}
\begin{tabular}{c l l}
  & {\bf Relation} &{\bf Error} \\ \hline \\
1. & $Z=cA$         & $\sigma(Z)=|c|\sigma(A)$  \\
2. & $Z=A+B+C+\cdots$ & $\sigma(Z)=\sqrt{(\sigma(A))^2 + (\sigma(B))^2
  +(\sigma(C))^2 + \cdots}$ \\
3. & $Z=A-B-C-\cdots$ & $\sigma(Z)=\sqrt{(\sigma(A))^2 + (\sigma(B))^2
  +(\sigma(C))^2 + \cdots}$   \\
4. & $Z=A \times B \times C \times \cdots$
  & $\frac{\sigma(Z)}{|Z|} = \sqrt{\left(\frac{\sigma(A)}{|A|}\right)^2 + \left(\frac{\sigma(B)}{|B|}\right)^2
  + \left(\frac{\sigma(C)}{|C|}\right)^2 + \cdots}$ \\
5. & $Z=\frac{A}{B}$ & $\frac{\sigma(Z)}{|Z|}
  =\sqrt{\left(\frac{\sigma(A)}{|A|}\right)^2+\left(\frac{\sigma(B)}{|B|}\right)^2}$ \\
6. & $Z=A^{b}$      & $\frac{\sigma(Z)}{|Z|}
  =\sqrt{\left(|b|~\frac{\sigma(A)}{|A|}\right)^2}$  \\
\end{tabular}
\end{center}
%%\label{table:errorrules}
%%\end{table}
Note that in the table, error terms are {\em always\/} added even when the relation is a subtraction.
\begin{itemize}
%%\vspace{-15mm}
\addtolength{\itemsep}{-0.5\baselineskip}
\item
$a, b, c, \ldots, z$ represent constants
\item
$A, B, C, \ldots, Z$ represent measured or calculated quantities
\item
$\sigma(A), \sigma(B), \sigma(C), \ldots, \sigma(Z)$ represent the errors in
$A, B, C, \ldots, Z$, respectively.
\end{itemize}

\section*{How to derive an error equation}
Let's use the change of variable method to determine the error equation for the
following expression: 

\begin{equation}
  y = \frac{M}{m}\sqrt{~0.5~kx} \\
  \label{eq:original}
\end{equation}

\begin{itemize}
\item
Begin by rewriting Equation~\ref{eq:original} as a product of terms:
\end{itemize}
\begin{eqnarray}
  y & = & M~*~m^{-1}~*~ \left[~0.5~*~k~*~x~*\right]~^{1/2} \\ 
    & = & M~*~m^{-1}~*~0.5^{1/2}~*~k^{1/2}~*~x^{1/2}~ 
  \label{eq:productofterms}
\end{eqnarray}

\begin{itemize}
\item
Assign to each term in Equation~\ref{eq:productofterms} a new variable name 
$A,B,C,\dots$~, then express $v$ in terms of these new variables, 
\end{itemize}
\begin{equation}
  y = A~*~B~*~C~*~D~*~E
\label{eq:product}
\end{equation}

\begin{itemize}
\item
With $\sigma(y)$ representing the error or uncertainty in the magnitude of $y$,
the error expression for $y$ is easily obtained by applying Rule 4 to the 
product of terms Equation~\ref{eq:product}: 
\end{itemize}

\begin{equation}
  \frac{\sigma(y)}{|y|} = \sqrt{ \left(\frac{\sigma(A)}{|A|}\right)^2~+~\left(\frac{\sigma(B)}{|B|}\right)^2~+~
  \left(\frac{\sigma(C)}{|C|}\right)^2~+~\left(\frac{\sigma(D)}{|D|}\right)^2~+~\left(\frac{\sigma(E)}{|E|}\right)^2}
  \label{eq:productoferrors}
\end{equation}

\begin{itemize}
\item
Select from the table of error rules an appropriate error expression 
for each of these new variables as shown below. Note that $F$ requires further
simplification since there are two terms under the square root, so we equate these
to a variable $G$: 
\end{itemize}
\begin{center}
\begin{tabular}{lll}
 $A = M$, & $\sigma(A) = \sigma(M)$ & \mbox{Rule 1} \\ 
 $B = m^{-1}$, & $\frac{\sigma(B)}{|B|} 
    = |-1|~\frac{\sigma(m)}{|m|} = \frac{\sigma(m)}{|m|}$ & \mbox{Rule 6} \\
 $C = 0.5^{1/2}$, & $\frac{\sigma(C)}{|C|} 
    = \left|\frac{1}{2}\right|~\frac{\sigma(0.5)}{|0.5|} 
    = 0$ & \mbox{since $\sigma(0.5)=0$} \\ 
 $D = k^{1/2}$, & $\frac{\sigma(D)}{|D|} 
    = \left|\frac{1}{2}\right|~\frac{\sigma(k)}{|k|} = \frac{\sigma(k)}{2|k|}$ 
      & \mbox{Rule 6}\\
 $E = x^{1/2}$, & $\frac{\sigma(E)}{|E|} 
    = \left|\frac{1}{2}\right|~\frac{\sigma(x)}{|x|} 
    = \frac{\sigma(x)}{2|x|}$ & \mbox{Rule 6} \\
\end{tabular}
\end{center}

\begin{itemize}
\item
Finally, replace the error terms into the original error 
Equation~\ref{eq:productoferrors}, simplify and solve for $\sigma(y)$ 
by multiplying both sides of the equation with $y$:
\end{itemize}

\begin{equation}
  \sigma(y) = |y|\sqrt{\left(\frac{\sigma(M)}{|M|}\right)^2 + \left(\frac{\sigma(m)}{|m|}\right)^2 + 
  \left(\frac{\sigma(k)}{2|k|}\right)^2 + \left(\frac{\sigma(x)}{2|x|}\right)^2} 
\end{equation}







%\end{document}



  % need work
%\include{./right_triangle/trigbook-chapter1} %god
%\include{./general_triangle/trigbook-chapter2} %good
%\input{./radian/trigbook-chapter4.tex} %good
%\import{./graphing_and_inverse/}{trigbook-chapter5} %good
%\import{./additional_topics/}{trigbook-chapter6} %good
% %\import{./AppExt/}{appext}

%\import{./vector/}{calc3book-chapter1_B} % need to correct some references

%\import{./topic/}{topic}
%%%%%%%%%%%%%%%%%%%%%%%%%%%%%%%%%%%%%%%%%

\newpage
%Put the bibliography here
\include{trigbook-biblio}
%Put appendices here
\fancyhf{}
\fancyhead[LE]{\tikzstyle{headbox} = [fill=headercolor,line width=0pt,rectangle,rounded corners]
 \begin{tikzpicture}
  \node [headbox] (box){%
   \begin{minipage}{\fminilength}%
     \sffamily{\bfseries\thepage\qquad \enskip\leftmark}
    \end{minipage}%
  };
 \end{tikzpicture}}
\fancyhead[LO]{\tikzstyle{headbox} = [fill=headercolor,line width=0pt,rectangle,rounded corners]
 \begin{tikzpicture}
  \node [headbox] (box){%
   \begin{minipage}{\fminilength}%
     \sffamily {}\hfill{\bfseries{\leftmark\qquad\thepage}}
    \end{minipage}%
  };
 \end{tikzpicture}}

%\include{./appendix/trigbook-appendixa}
%\include{./appendix/trigbook-appendixb}
%\include{conversion_factor}
%\include{conversion}
\newpage
%Put the GNU PDL here
\include{gnupdl}
%Put the history section here

\backmatter
\clearpage
\phantomsection
%Put the index here
\printindex

\end{document}
